% =============================================================
% main.tex — 図でつながる ITパスポート
% B5 / LuaLaTeX / jlreq
% =============================================================
\documentclass[
  paper=b5j,
  fontsize=10pt,
  book,
  head_space=25mm,
  foot_space=25mm,
  gutter=22mm,
  fore-edge=18mm,
  twoside,
  openany,
]{jlreq}

% ── スタイル読み込み ──
% =============================================================
% colors.tex — カラーパレット一元管理
% =============================================================
\usepackage{xcolor}

% ── プライマリ ──
\definecolor{Navy}{HTML}{1F3A5F}
\definecolor{Teal}{HTML}{0E7C86}
\definecolor{SubBlue}{HTML}{3C5A80}
\definecolor{Ink}{HTML}{1C2430}
\definecolor{Muted}{HTML}{7A8698}
\definecolor{RuleLine}{HTML}{C8D0DA}
\definecolor{HairLine}{HTML}{E4E9EF}
\definecolor{PageBG}{HTML}{FFFFFF}

% ── 難易度バッジ ──
\definecolor{BadgeKiso}{HTML}{2C7A46}
\definecolor{BadgeHyojun}{HTML}{2B67A0}
\definecolor{BadgeKyoka}{HTML}{C8662A}
\definecolor{BadgeHatten}{HTML}{8B2252}

% ── PART / 章扉 ──
\definecolor{ChapNumBG}{HTML}{E8EDF5}
\definecolor{TopBar}{HTML}{1F3A5F}
\definecolor{AccentBar}{HTML}{0E7C86}
\definecolor{LeadBG}{HTML}{F7F9FC}

% ── 例題 (reidai) ──
\definecolor{ReidaiBar}{HTML}{1F3A5F}
\definecolor{ReidaiBG}{HTML}{F2F5FA}

% ── 指針 (shishin) ──
\definecolor{ShishinBar}{HTML}{0E7C86}
\definecolor{ShishinBG}{HTML}{EAF5F4}

% ── 解答 (kaitou) ──
\definecolor{KaitouBar}{HTML}{2C7A46}
\definecolor{KaitouBG}{HTML}{F0F9F2}

% ── 見抜き方 (minuki) ──
\definecolor{MinukiBar}{HTML}{5B4A8B}
\definecolor{MinukiBG}{HTML}{F1ECF5}

% ── 定義 / 要点 (teigi) ──
\definecolor{TeigiBar}{HTML}{0E7C86}
\definecolor{TeigiBG}{HTML}{EAF5F4}

% ── つまずき (pitfall) ──
\definecolor{PitBar}{HTML}{C8662A}
\definecolor{PitBG}{HTML}{FBEFE3}

% ── 比較 (hikaku) ──
\definecolor{HikakuBar}{HTML}{3C5A80}
\definecolor{HikakuBG}{HTML}{ECF1F7}

% ── 解説 (kaisetsu) ──
\definecolor{KaiBar}{HTML}{A8934F}
\definecolor{KaiBG}{HTML}{FBF6E3}

% ── よくあるミス (yokumiss) ──
\definecolor{MissBar}{HTML}{B6551F}
\definecolor{MissBG}{HTML}{FFF1E6}

% ── 到達チェック ──
\definecolor{CheckGreen}{HTML}{2C7A46}

% ── 問題 ──
\definecolor{QuizBar}{HTML}{1F3A5F}
\definecolor{OrigBadge}{HTML}{245C4F}
\definecolor{OfficialBadge}{HTML}{7B5A1E}

% =============================================================
% layout.tex — B5 紙面レイアウト / ヘッダ・フッタ / 見出し / 目次
% =============================================================

% ── 日本語 ──
\usepackage{luatexja}
\usepackage{luatexja-fontspec}

% ── ページ余白 ──
\usepackage[b5paper,
  top=25mm, bottom=25mm,
  inner=22mm, outer=18mm,
  headheight=14pt, headsep=10mm, footskip=14mm
]{geometry}

% ── パッケージ群 ──
\usepackage{fancyhdr}
\usepackage{enumitem}
\usepackage{parskip}
\usepackage{graphicx}
\usepackage{adjustbox}
\usepackage{array}
\usepackage{booktabs}
\usepackage{longtable}
\usepackage{tabularx}
\usepackage{colortbl}
\usepackage{multicol}
\usepackage{tikz}
\usetikzlibrary{positioning,arrows.meta,shapes.geometric,fit,backgrounds,calc,shadows.blur}
\usepackage[most]{tcolorbox}
\tcbuselibrary{skins,breakable,xparse}
\usepackage{hyperref}
\usepackage{amssymb}
\usepackage{titletoc}
\usepackage{etoolbox}
\usepackage{xparse}
\usepackage{pgffor}

% ── ハイパーリンク ──
\hypersetup{
  colorlinks=true,
  linkcolor=Navy,
  urlcolor=Teal,
  bookmarksnumbered=true,
  pdftitle={図でつながる ITパスポート},
}

% =============================================================
% 章扉デザイン: jlreq では \chapter をそのまま使い、
% 装飾は \chaptermeta 等のコマンドで付加する方針とする。
% =============================================================

% =============================================================
% section / subsection デザイン（jlreq 互換）
% =============================================================

% =============================================================
% 目次デザイン
% =============================================================
\setcounter{tocdepth}{1}
\renewcommand{\contentsname}{目次}

\titlecontents{chapter}[0pt]
  {\addvspace{10pt}}
  {\large\bfseries\color{Navy}\makebox[2.2em][l]{\thecontentslabel}}
  {\large\bfseries\color{Navy}}
  {\normalfont\color{Muted}\ \titlerule*[0.7pc]{\footnotesize.}\ \contentspage}
  [\vspace{1pt}{\color{HairLine}\rule{\linewidth}{0.3pt}}]

\titlecontents{section}[2em]
  {\addvspace{2pt}\small}
  {\color{SubBlue}\makebox[2em][l]{\thecontentslabel}}
  {}
  {\normalfont\footnotesize\color{Muted}\ \titlerule*[0.7pc]{\tiny.}\ \contentspage}

% =============================================================
% ヘッダ / フッタ
% =============================================================
\pagestyle{fancy}
\fancyhf{}
\fancyhead[LE]{\small\color{Muted}\sffamily 図でつながる ITパスポート}
\fancyhead[RO]{\small\color{SubBlue}\nouppercase\leftmark}
\fancyfoot[C]{\small\color{Muted}\thepage}
\renewcommand{\headrulewidth}{0pt}
\renewcommand{\footrulewidth}{0pt}
\renewcommand{\headrule}{%
  \hbox to\headwidth{\color{HairLine}\leaders\hrule height 0.3pt\hfill}%
}
\makeatletter
\def\chaptermark#1{\markboth{第\,\thechapter\,章\quad #1}{}}
\makeatother

% ── リスト / 段落 ──
\setlist[itemize]{leftmargin=1.4em, itemsep=2pt, topsep=4pt}
\setlist[enumerate]{leftmargin=1.6em, itemsep=2pt, topsep=4pt}
\setlength{\parindent}{1em}
\setlength{\parskip}{0.4em plus 0.1em minus 0.05em}
\renewcommand{\baselinestretch}{1.08}

% =============================================================
% environments.tex — 全 tcolorbox 環境
% =============================================================

% ── 共通バー幅 ──
\newcommand{\BarW}{3.6pt}
\newcommand{\ArcR}{1.5pt}

% ── 汎用ラベル付き左バーボックス ──
\NewDocumentCommand{\MakeActorBox}{m m m m m}{%
  % #1=env name, #2=default label, #3=bar color, #4=bg color, #5=title bg
  \newtcolorbox{#1}[1][#2]{%
    enhanced, breakable,
    colback=#4, colframe=#4, boxrule=0pt,
    leftrule=\BarW,
    colbacktitle=#5, coltitle=white,
    fonttitle=\sffamily\footnotesize\bfseries,
    title=\,##1\,,
    attach boxed title to top left={xshift=-\BarW, yshift=-6pt},
    boxed title style={colframe=#5, boxrule=0pt, sharp corners, arc=0pt,
      left=6pt, right=6pt, top=2pt, bottom=2pt},
    overlay unbroken and first={
      \draw[#3, line width=\BarW] (frame.north west) -- (frame.south west);
    },
    arc=\ArcR, outer arc=0pt, sharp corners=northwest,
    left=12pt, right=10pt, top=14pt, bottom=8pt,
    fontupper=\color{Ink}\small,
    before skip=10pt, after skip=10pt,
  }%
}

% ── 定義/要点 ──
\MakeActorBox{teigi}{要点}{TeigiBar}{TeigiBG}{TeigiBar}

% ── つまずき ──
\MakeActorBox{pitfall}{つまずき}{PitBar}{PitBG}{PitBar}

% ── 比較 ──
\MakeActorBox{hikaku}{比較}{HikakuBar}{HikakuBG}{HikakuBar}

% ── 見抜き方 ──
\MakeActorBox{minuki}{見抜き方}{MinukiBar}{MinukiBG}{MinukiBar}

% ── よくあるミス ──
\MakeActorBox{yokumiss}{よくあるミス}{MissBar}{MissBG}{MissBar}

% ── 解説 ──
\MakeActorBox{kaisetsubox}{解答・解説}{KaiBar}{KaiBG}{KaiBar}

% ── アドバイス ──
\MakeActorBox{advice}{学習アドバイス}{Muted}{LeadBG}{Muted}

% ── 本試験パターン ──
\MakeActorBox{honshiken}{本試験}{MinukiBar}{MinukiBG}{MinukiBar}

% =============================================================
% 例題 (reidai): 紺ヘッダーバー＋問題
% =============================================================
\newcounter{reidaictr}
\renewcommand{\thereidaictr}{\arabic{reidaictr}}

\NewDocumentEnvironment{reidai}{m}{%
  \refstepcounter{reidaictr}%
  \par\addvspace{14pt}%
  \begin{tcolorbox}[
    enhanced, breakable,
    colback=ReidaiBG, colframe=ReidaiBar,
    boxrule=0pt, toprule=3pt,
    arc=0pt, outer arc=0pt,
    left=12pt, right=12pt, top=10pt, bottom=8pt,
    before skip=10pt, after skip=4pt,
    title={\sffamily\small\bfseries 例題 \thereidaictr\quad #1},
    colbacktitle=ReidaiBar, coltitle=white,
    attach boxed title to top left={yshift=0pt, xshift=0pt},
    boxed title style={colframe=ReidaiBar, boxrule=0pt, arc=0pt,
      left=10pt, right=10pt, top=4pt, bottom=4pt},
  ]
  \color{Ink}\small
}{%
  \end{tcolorbox}
  \par\addvspace{4pt}%
}

% ── 指針 ──
\NewDocumentEnvironment{shishin}{}{%
  \begin{tcolorbox}[
    enhanced, breakable,
    colback=ShishinBG, colframe=ShishinBG, boxrule=0pt,
    leftrule=\BarW,
    overlay unbroken and first={
      \draw[ShishinBar, line width=\BarW] (frame.north west) -- (frame.south west);
    },
    arc=\ArcR, outer arc=0pt,
    left=12pt, right=10pt, top=10pt, bottom=8pt,
    before skip=4pt, after skip=4pt,
    fontupper=\color{Ink}\small,
    title={\sffamily\footnotesize\bfseries\color{ShishinBar} 指針},
    coltitle=ShishinBar, colbacktitle=ShishinBG,
    attach boxed title to top left={xshift=8pt, yshift=-2pt},
    boxed title style={colframe=ShishinBG, boxrule=0pt, arc=0pt,
      left=0pt, right=0pt, top=0pt, bottom=0pt},
  ]
}{%
  \end{tcolorbox}
}

% ── 解答（緑左バー）──
\NewDocumentEnvironment{kaitou}{}{%
  \par\addvspace{4pt}%
  \begin{tcolorbox}[
    enhanced, breakable,
    colback=KaitouBG, colframe=KaitouBG, boxrule=0pt,
    leftrule=\BarW,
    colbacktitle=KaitouBar, coltitle=white,
    fonttitle=\sffamily\footnotesize\bfseries,
    title=\,解答\,,
    attach boxed title to top left={xshift=-\BarW, yshift=-6pt},
    boxed title style={colframe=KaitouBar, boxrule=0pt, sharp corners, arc=0pt,
      left=6pt, right=6pt, top=2pt, bottom=2pt},
    overlay unbroken and first={
      \draw[KaitouBar, line width=\BarW] (frame.north west) -- (frame.south west);
    },
    arc=\ArcR, outer arc=0pt,
    left=12pt, right=10pt, top=14pt, bottom=8pt,
    fontupper=\color{Ink}\small,
    before skip=4pt, after skip=10pt,
  ]
}{%
  \end{tcolorbox}
}

% =============================================================
% 4択問題 (itpmon): ア/イ/ウ/エ 本試験準拠
% =============================================================
\newcounter{itpmonctr}
\renewcommand{\theitpmonctr}{\arabic{itpmonctr}}

% #1 = key-value options, #2 = problem text
\NewDocumentEnvironment{itpmon}{O{} m}{%
  \refstepcounter{itpmonctr}%
  \par\addvspace{8pt}%
  \begin{tcolorbox}[
    enhanced, breakable,
    colback=white, colframe=HairLine,
    boxrule=0.5pt, arc=0pt, outer arc=0pt,
    left=10pt, right=10pt, top=10pt, bottom=8pt,
    before skip=8pt, after skip=4pt,
  ]
    \noindent
    \tikz[baseline=(q.base)]{
      \node[fill=QuizBar, text=white, inner xsep=7pt, inner ysep=2pt,
            font=\sffamily\scriptsize\bfseries, rounded corners=1pt] (q)
        {Q\arabic{itpmonctr}};
    }%
    \hspace{4pt}%
    {\small\color{Muted}#1}%
    \par\vspace{5pt}%
    \textbf{\color{Navy}#2}\par
    \vspace{4pt}%
}{%
  \end{tcolorbox}
}

% 選択肢マクロ: 4つの選択肢を表示
\NewDocumentCommand{\sentakushi}{m m m m}{%
  \begin{enumerate}[label=\textbf{\color{SubBlue}\ifcase\value{enumi}\or ア\or イ\or ウ\or エ\fi.},
                   leftmargin=2.2em, itemsep=3pt, topsep=4pt]
    \item #1
    \item #2
    \item #3
    \item #4
  \end{enumerate}
}

% 解説マクロ: 正解＋解説
% #1 = 正解 (ア/イ/ウ/エ)
\NewDocumentEnvironment{kaisetsu}{m}{%
  \begin{kaisetsubox}
    \textbf{\color{KaiBar}正解}\quad\textbf{#1}\par
    \vspace{3pt}%
}{%
  \end{kaisetsubox}
}

% =============================================================
% 到達確認チェックリスト
% =============================================================
\NewDocumentEnvironment{tatsaku}{O{到達確認チェックリスト}}{%
  \par\addvspace{12pt}%
  \begin{tcolorbox}[
    enhanced, breakable,
    colback=white, colframe=Teal,
    boxrule=0.6pt, arc=2pt, leftrule=4pt,
    left=12pt, right=10pt, top=10pt, bottom=10pt,
    before skip=12pt, after skip=12pt,
    title={\sffamily\small\bfseries #1},
    colbacktitle=Teal, coltitle=white,
    attach boxed title to top left={xshift=-4pt, yshift=-6pt},
    boxed title style={colframe=Teal, boxrule=0pt, sharp corners, arc=0pt,
      left=8pt, right=8pt, top=2pt, bottom=2pt},
  ]
  \small\color{Ink}%
  \begin{itemize}[leftmargin=1.8em, itemsep=3pt, topsep=2pt,
    label={\color{Teal}$\square$}]
}{%
  \end{itemize}
  \end{tcolorbox}
}

% =============================================================
% タイムアタック
% =============================================================
\NewDocumentEnvironment{timeattack}{m m m}{%
  \par\addvspace{14pt}%
  \noindent
  \tikz[baseline=(t.base)]{\node[fill=Navy, text=white, inner xsep=8pt, inner ysep=3pt,
    font=\sffamily\scriptsize\bfseries, rounded corners=1pt] (t) {TIME ATTACK};}%
  \hspace{4pt}%
  \tikz[baseline=(u.base)]{\node[fill=Teal, text=white, inner xsep=6pt, inner ysep=2pt,
    font=\sffamily\scriptsize\bfseries, rounded corners=1pt] (u) {#2};}%
  \hspace{4pt}%
  \tikz[baseline=(v.base)]{\node[fill=BadgeHyojun, text=white, inner xsep=6pt, inner ysep=2pt,
    font=\sffamily\scriptsize\bfseries, rounded corners=1pt] (v) {目標 #3};}%
  \par\vspace{6pt}%
  {\Large\bfseries\color{Navy}#1\par}%
  \vspace{3pt}%
  {\color{HairLine}\rule{\linewidth}{0.4pt}}\par
  \vspace{6pt}%
}{%
  \par\vspace{6pt}%
  {\color{HairLine}\rule{\linewidth}{0.4pt}}\par
  \vspace{10pt}%
}

% =============================================================
% 模試形式セット
% =============================================================
\NewDocumentEnvironment{mogishi}{m m m}{%
  \clearpage
  \noindent
  \tikz[baseline=(m.base)]{\node[fill=Navy, text=white, inner xsep=10pt, inner ysep=4pt,
    font=\sffamily\small\bfseries, rounded corners=1pt] (m) {模試形式};}%
  \hspace{4pt}%
  \tikz[baseline=(n.base)]{\node[fill=Teal, text=white, inner xsep=8pt, inner ysep=3pt,
    font=\sffamily\scriptsize\bfseries, rounded corners=1pt] (n) {#2};}%
  \hspace{4pt}%
  \tikz[baseline=(o.base)]{\node[fill=BadgeHyojun, text=white, inner xsep=8pt, inner ysep=3pt,
    font=\sffamily\scriptsize\bfseries, rounded corners=1pt] (o) {目標 #3};}%
  \par\vspace{8pt}%
  {\color{TopBar}\rule{0.4\linewidth}{4pt}}%
  {\color{AccentBar}\rule{0.2\linewidth}{4pt}}\par
  \vspace{8pt}%
  {\LARGE\bfseries\color{Navy}#1\par}%
  \vspace{6pt}%
  {\color{HairLine}\rule{\linewidth}{0.4pt}}\par
  \vspace{8pt}%
}{%
  \par\vspace{10pt}%
  {\color{TopBar}\rule{\linewidth}{1pt}}\par
  \vspace{12pt}%
}

% =============================================================
% まとめボックス
% =============================================================
\MakeActorBox{matome}{この章のまとめ}{Navy}{LeadBG}{Navy}

% =============================================================
% commands.tex — ユーティリティコマンド
% =============================================================

% ── 汎用バッジ ──
\NewDocumentCommand{\badge}{m m}{%
  \tikz[baseline=(b.base)]{
    \node[fill=#1, text=white, inner xsep=6pt, inner ysep=2pt,
          font=\sffamily\scriptsize\bfseries, rounded corners=1pt] (b) {#2};
  }%
}

% ── 難易度バッジ ──
\NewDocumentCommand{\diffbadge}{m}{%
  \ifstrequal{#1}{基}{\badge{BadgeKiso}{基\,基礎}}{%
  \ifstrequal{#1}{標}{\badge{BadgeHyojun}{標\,標準}}{%
  \ifstrequal{#1}{強}{\badge{BadgeKyoka}{強\,強化}}{%
  \ifstrequal{#1}{発}{\badge{BadgeHatten}{発\,発展}}{%
    \badge{BadgeHyojun}{#1}%
  }}}}%
}

% ── 出典バッジ ──
\NewDocumentCommand{\sourcebadge}{m}{%
  \ifstrequal{#1}{original}{\badge{OrigBadge}{オリジナル}}{%
    \badge{OfficialBadge}{#1}%
  }%
}

% ── 強調タグ ──
\newcommand{\termtag}[1]{\textcolor{Teal}{\textbf{#1}}}

% ── 比較表テンプレ ──
\NewDocumentCommand{\comptable}{m m}{%
  \begin{center}
  \small
  \renewcommand{\arraystretch}{1.2}%
  \arrayrulecolor{HairLine}%
  \begin{tabularx}{0.96\linewidth}{@{}>{\color{SubBlue}\bfseries\sffamily\footnotesize}l X X@{}}
    \arrayrulecolor{HikakuBar}\toprule
    #1 \\
    \arrayrulecolor{HairLine}\midrule
    #2 \\
    \arrayrulecolor{HikakuBar}\bottomrule
  \end{tabularx}
  \arrayrulecolor{black}%
  \end{center}%
}

% ── PART 区切りページ ──
\NewDocumentCommand{\partpage}{m m}{%
  \clearpage
  \thispagestyle{empty}%
  \null\vfill
  \begin{center}
    {\color{Navy!15}\rule{0.6\linewidth}{0.4pt}}\par
    \vspace{16mm}%
    {\sffamily\Large\color{Muted}PART\enspace #1}\par
    \vspace{10mm}%
    {\fontsize{28pt}{34pt}\selectfont\bfseries\color{Navy}#2\par}%
    \vspace{10mm}%
    {\color{AccentBar}\rule{0.2\linewidth}{2.5pt}}\par
  \end{center}
  \vfill\null
  \clearpage
}

% ── 章扉メタ情報 ──
% 難易度の星
\newcommand{\diffstars}[1]{%
  \foreach \i in {1,...,4}{%
    \ifnum\i>#1\relax
      {\color{HairLine}$\star$}%
    \else
      {\color{Teal}$\bigstar$}%
    \fi
  }%
}

% #1=スキルリスト, #2=目安時間, #3=難易度(1-4), #4=前提知識
\NewDocumentCommand{\chaptermeta}{m m m m}{%
  \begin{tcolorbox}[
    enhanced, breakable,
    colback=LeadBG, colframe=HairLine,
    boxrule=0.4pt, arc=2pt,
    left=14pt, right=14pt, top=12pt, bottom=12pt,
    before skip=4pt, after skip=16pt,
  ]
    {\sffamily\small\bfseries\color{Teal}この章で身につく力}\par\vspace{4pt}%
    {\small\color{Ink}%
      \begin{itemize}[leftmargin=1.2em, itemsep=1pt, topsep=2pt]
        #1
      \end{itemize}
    }%
    \vspace{6pt}%
    {\color{HairLine}\rule{\linewidth}{0.3pt}}\par\vspace{6pt}%
    \noindent
    \begin{minipage}[t]{0.30\linewidth}
      {\sffamily\scriptsize\color{Muted}目安時間}\par
      {\large\bfseries\color{Navy}#2}
    \end{minipage}%
    \begin{minipage}[t]{0.30\linewidth}
      {\sffamily\scriptsize\color{Muted}難易度}\par
      \diffstars{#3}
    \end{minipage}%
    \begin{minipage}[t]{0.40\linewidth}
      {\sffamily\scriptsize\color{Muted}前提知識}\par
      {\footnotesize\color{Ink}#4}
    \end{minipage}
  \end{tcolorbox}
}

% ── 次章への橋渡し ──
\newcommand{\nextchapter}[1]{%
  \par\addvspace{12pt}%
  {\color{AccentBar}\rule{0.3\linewidth}{2pt}}\par\vspace{4pt}%
  {\small\bfseries\color{Navy}次の章では：}%
  {\small\color{Ink}#1\par}%
  \vspace{8pt}%
}

% ── セクション区切りバー ──
\newcommand{\secbar}{%
  \par\vspace{8pt}%
  {\color{TopBar}\rule{0.38\linewidth}{4pt}}%
  {\color{AccentBar}\rule{0.18\linewidth}{4pt}}\par
  \vspace{8pt}%
}

% ── 章末演習見出し ──
\newcommand{\shoumatsuhead}[1]{%
  \par\addvspace{14pt}%
  \noindent\badge{Navy}{#1}%
  \hspace{6pt}%
  {\color{HairLine}\leaders\hrule height 0.4pt\hfill\null}\par
  \vspace{6pt}%
}


\begin{document}

% =============================================================
% 表紙
% =============================================================
\begin{titlepage}
\thispagestyle{empty}
\noindent
{\color{TopBar}\rule{\linewidth}{6pt}}\par
\vspace{2pt}
{\color{AccentBar}\rule{0.4\linewidth}{2pt}}\par
\vspace{30mm}

\noindent
\badge{Navy}{ITパスポート\ 問題集}
\par\vspace{14mm}

{\raggedright
  \fontsize{30pt}{38pt}\selectfont\bfseries\color{Navy}
  図でつながる\par
  \vspace{4pt}
  ITパスポート\par
}
\vspace{8mm}
{\color{AccentBar}\rule{0.25\linewidth}{3pt}}\par
\vspace{6mm}
{\large\color{SubBlue}
  ―― 過去問と図解で理解し、\\[2pt]
  \hspace*{1.2em}その場で解ける問題集\par
}
\vfill

\noindent
\begin{minipage}[b]{0.6\linewidth}
  \small\color{Muted}
  図で理解する。すぐ解く。解説でひっかけを潰す。\\[4pt]
  本書は暗記ではなく\textbf{\color{Navy}見分け方の理解}で\\
  合格を目指す問題集です。\\[4pt]
  令和6年度・令和7年度 公開問題収録
\end{minipage}%
\hfill
\begin{minipage}[b]{0.35\linewidth}
  \raggedleft
  {\footnotesize\sffamily\color{Muted}2026年版}
\end{minipage}
\vspace{6mm}

{\color{AccentBar}\rule{0.4\linewidth}{2pt}}\par
\vspace{2pt}
{\color{TopBar}\rule{\linewidth}{6pt}}
\end{titlepage}

% =============================================================
% 目次
% =============================================================
\pagenumbering{roman}
\pagestyle{empty}
\cleardoublepage

\vspace*{8mm}
\noindent\badge{Navy}{CONTENTS}\par
\vspace{4pt}
{\color{TopBar}\rule{0.24\linewidth}{3pt}}%
{\color{AccentBar}\rule{0.14\linewidth}{3pt}}\par
\vspace{6pt}
{\raggedright\Huge\bfseries\color{Navy}目\,次\par}
\vspace{2pt}
{\color{HairLine}\rule{\linewidth}{0.4pt}}\par
\vspace{8mm}

\makeatletter
\renewcommand{\contentsname}{}%
\@starttoc{toc}
\makeatother

\cleardoublepage
\pagestyle{fancy}

% =============================================================
% 本文
% =============================================================
\pagenumbering{arabic}
\setcounter{page}{1}

% =============================================================
% 00_howto.tex — 本書の使い方
% =============================================================
\chapter*{この本の使い方}
\addcontentsline{toc}{chapter}{この本の使い方}

\begin{teigi}[ITパスポート試験とは]
ITパスポート試験は、ITに関する基礎的な知識を証明する\textbf{国家試験}です。

\begin{itemize}
  \item \textbf{試験時間}：120分
  \item \textbf{問題数}：100問（全問必須・四肢択一）
  \item \textbf{合格基準}：総合600点以上 ＋ 各分野300点以上
  \item \textbf{出題分野}：ストラテジ系（問1〜35）、マネジメント系（問36〜55）、テクノロジ系（問56〜100）
\end{itemize}
\end{teigi}

\section*{本書の構成}

本書は\textbf{4つのパート}で構成されています。

\begin{hikaku}[4パート構成]
\begin{itemize}
  \item \textbf{Part I}（第1〜2章）：ストラテジ系——企業活動・経営戦略・会計・法務
  \item \textbf{Part II}（第3〜4章）：マネジメント系——開発・PM・サービス・監査
  \item \textbf{Part III}（第5〜7章）：テクノロジ系——基礎理論・NW・DB・セキュリティ・AI
  \item \textbf{Part IV}（第8〜9章）：総仕上げ——横断整理・模擬試験3セット
\end{itemize}
\end{hikaku}

\section*{各章の読み方}

各章は次の流れで構成されています。

\begin{enumerate}
  \item \textbf{この章で身につく力}——学習目標と目安時間
  \item \textbf{要点}——見抜き方のルール
  \item \textbf{例題}——ステップ付きの解法（指針→解答→見抜き方）
  \item \textbf{過去問ドリル}——R6/R7公開問題の演習
  \item \textbf{タイムアタック}——制限時間つきの実戦演習
  \item \textbf{よくあるミス}——ひっかけパターンの対策
  \item \textbf{到達確認チェックリスト}——自己診断
\end{enumerate}

\section*{難易度の表示}

問題には次の難易度が付いています。

\begin{itemize}
  \item \diffbadge{基}\quad 基礎——1つの知識で直接解ける問題
  \item \diffbadge{標}\quad 標準——正しい知識の選択が必要な問題
  \item \diffbadge{強}\quad 強化——複数知識の組合せや、ひっかけを含む問題
  \item \diffbadge{発}\quad 発展——計算や読解力が必要な問題
\end{itemize}

まずは\diffbadge{基}と\diffbadge{標}を確実にしてから、\diffbadge{強}に進みましょう。

\section*{出典について}

\begin{honshiken}[出典表記]
本書に収録する過去問は、独立行政法人情報処理推進機構（IPA）が公表している公開問題に基づきます。
問題文は公開問題の記述を尊重し、学習目的の範囲で引用・参照しています。
各問の冒頭に出典年度を\sourcebadge{R7}のようにバッジで示します。
\end{honshiken}


% ── PART I: ストラテジ系 ──
\partpage{I}{ストラテジ系}
% =============================================================
% 01_enterprise.tex — 企業活動・経営戦略
% =============================================================
\chapter{企業活動と経営戦略}

\chaptermeta{%
  \item 経営理念→方針→戦略→戦術の4段階を説明できる
  \item CxO の各役職の違いを即答できる
  \item SWOT/VRIO/PPM 等のフレームワークを使い分けられる
  \item DX/PoC/リーンスタートアップの関係を整理できる
  \item 新事業・起業の用語を正しく識別できる
}{75分}{2}{なし（本書の出発点）}

Iパス最大の出題分野が「ストラテジ系」です。
本章は\textbf{企業活動 → 経営戦略フレームワーク → IT活用}の順に進みます。
すべて\textbf{「経営者の意思決定」}という1本の軸で貫かれています。

% =============================================================
\section{企業活動と組織}
% =============================================================

\begin{pitfall}[「経営理念」と「経営戦略」がごちゃごちゃ]
抽象から具体への4段階を固定しましょう。
\begin{itemize}
  \item \termtag{経営理念}：企業が存在する根本の目的（なぜ）
  \item \termtag{経営方針}：理念を実現する大きな方向性（どちらへ）
  \item \termtag{経営戦略}：方針を実現する具体的な道筋（どうやって）
  \item \termtag{戦術}：戦略を実現する現場の施策（いま何を）
\end{itemize}
\end{pitfall}

\subsection{組織形態の整理}

\begin{hikaku}[代表的な組織形態]
\comptable{形態 & 特徴 & 強み／弱み}{
職能別組織 & 機能（営業・製造等）で分ける & 専門性↑／全体最適↓ \\
事業部制組織 & 事業ごとに独立採算 & 機動力↑／重複↑ \\
マトリックス組織 & 機能×プロジェクトの2軸 & 柔軟性↑／指揮系統が二重 \\
カンパニー制 & 事業を社内分社化 & 責任明確／小回り困難
}
\end{hikaku}

\subsection{経営層の役職}

\begin{teigi}[CxO：頭文字で役割がわかる]
\termtag{CEO}（最高経営責任者）／\termtag{COO}（執行）／\termtag{CFO}（財務）／
\termtag{CIO}（情報）／\termtag{CTO}（技術）／\termtag{CISO}（情報セキュリティ）。
「C = Chief、O = Officer、真ん中が担当分野」と覚える。
\end{teigi}

% --- 例題 ---
\begin{reidai}{CxOの識別}
技術戦略の策定や技術開発の推進といった技術経営に直接の責任をもつ役職はどれか。

\sentakushi{CEO}{CFO}{COO}{CTO}
\end{reidai}

\begin{shishin}
「技術経営」「技術開発の推進」というキーワードに注目する。
各CxOの真ん中のアルファベットが担当分野を示すので、「T = Technology」を探す。
\end{shishin}

\begin{kaitou}
\textbf{正解：エ（CTO）}

CTO = Chief \textbf{Technology} Officer（最高技術責任者）。
CEO は経営全体、CFO は財務、COO は業務執行の責任者。
\end{kaitou}

\begin{minuki}[見抜き方]
\begin{itemize}
  \item CxOは「真ん中のアルファベット」が守備範囲を示す
  \item 問題文のキーワード（技術/財務/情報/セキュリティ）と照合するだけ
  \item CIO（情報）とCTO（技術）の混同に注意——CIOは「情報\textbf{システム戦略}」、CTOは「\textbf{技術}開発」
\end{itemize}
\end{minuki}

% =============================================================
\section{経営戦略のフレームワーク}
% =============================================================

\begin{pitfall}[横文字フレームワークが多すぎる]
「\textbf{何を見るための道具か}」を決めて引き出しに収めれば怖くありません。
\begin{itemize}
  \item 自社を見る → SWOT、VRIO
  \item 市場を見る → 3C、5F（5フォース）
  \item 売り方を設計 → 4P／4C
  \item 投資配分を決める → PPM
  \item 成長方向を選ぶ → アンゾフの成長マトリクス
\end{itemize}
\end{pitfall}

\begin{hikaku}[戦略フレームワークの引き出し]
\comptable{フレーム & 観点 & 一言で}{
SWOT & 自社の健康診断 & 強み・弱み×機会・脅威 \\
VRIO & 強みの深堀り & 価値・希少性・模倣困難性・組織 \\
3C & 市場の三角測量 & 顧客・競合・自社 \\
5F & 業界の圧力 & 既存・新規・代替・売手・買手 \\
4P/4C & マーケティング & 売り手4視点／買い手4視点 \\
PPM & 投資配分 & 成長率×占有率の4象限 \\
アンゾフ & 成長方向 & 既存/新規 × 市場/製品
}
\end{hikaku}

\subsection{PPMの4象限}

\begin{teigi}[PPM（プロダクト・ポートフォリオ・マネジメント）]
\textbf{市場成長率}（縦軸）と\textbf{市場占有率}（横軸）で事業を4象限に分類する。
\begin{itemize}
  \item \termtag{花形（Star）}：成長率高・占有率高 → 投資継続
  \item \termtag{金のなる木（Cash Cow）}：成長率低・占有率高 → 収穫
  \item \termtag{問題児（Question Mark）}：成長率高・占有率低 → 選別投資
  \item \termtag{負け犬（Dog）}：成長率低・占有率低 → 撤退検討
\end{itemize}
\end{teigi}

% --- 例題 ---
\begin{reidai}{VRIO分析の識別}
投資の優先度などの経営の戦略を策定するために、経済価値、希少性、模倣困難性及び組織の四つの要素で評価することによって、自社のもつ資源を分析する手法として、最も適切なものはどれか。

\sentakushi{4P}{PPM}{SWOT分析}{VRIO分析}
\end{reidai}

\begin{shishin}
問題文に「経済価値、希少性、模倣困難性、組織」と4要素が列挙されている。
各フレームワークの「何を見る道具か」を思い出し、この4要素に一致するものを選ぶ。
\end{shishin}

\begin{kaitou}
\textbf{正解：エ（VRIO分析）}

VRIO = \textbf{V}alue（経済価値）・\textbf{R}areness（希少性）・\textbf{I}nimitability（模倣困難性）・\textbf{O}rganization（組織）。
4P はマーケティングミックス、PPM は成長率×占有率、SWOT は強み弱み×機会脅威。
\end{kaitou}

\begin{minuki}[見抜き方]
\begin{itemize}
  \item 問題文に「4つの要素」が列挙されていたら、各フレームワークの構成要素と照合する
  \item VRIO は「その強みは\textbf{本当に競争優位か}」を4つの問いで確認するツール
  \item SWOT は「内部×外部」の2軸、VRIO は「強みの質」を掘り下げる点が違う
\end{itemize}
\end{minuki}

% =============================================================
\section{新事業とIT活用}
% =============================================================

\begin{hikaku}[新規事業の関連用語]
\comptable{用語 & 要点 & 混同しやすい用語}{
PoC & 技術の実現性確認 & PoV（価値の実証）より前段階 \\
PoV & ビジネス価値の実証 & PoCの次のステップ \\
MVP & 検証用の最小限の製品 & 本番製品ではない \\
リーンスタートアップ & MVPで仮説検証を繰り返す手法 & DX（概念）とは階層が違う \\
DX & ITで経営を変革する概念 & 手法ではなく変革そのもの \\
ハッカソン & 短期集中でプロトタイプを作り競う & セキュリティキャンプとは別
}
\end{hikaku}

% --- 例題 ---
\begin{reidai}{PoCの識別}
新しい概念やアイディアの実証を目的とした、開発の前段階における検証を表す用語はどれか。

\sentakushi{CRM}{PoC}{RAS}{SLA}
\end{reidai}

\begin{shishin}
「概念の実証」「開発の前段階」というキーワードに注目。
Proof of Concept = 概念実証。CRMは顧客管理、RASは信頼性指標、SLAはサービス品質合意。
\end{shishin}

\begin{kaitou}
\textbf{正解：イ（PoC）}

PoC（Proof of Concept）は「概念実証」。新技術やアイデアが実現可能かを開発前に小規模に検証する活動。
\end{kaitou}

% =============================================================
\section{過去問ドリル}
% =============================================================

\shoumatsuhead{過去問ドリル：ストラテジ系}

\begin{itpmon}[\sourcebadge{R7}\enspace\diffbadge{標}]{%
ハッカソンに関する記述として、最も適切なものはどれか。}
\sentakushi
  {定められたルールの下、ある主題について、肯定派と否定派といった異なる立場に分かれて議論する。}
  {情報セキュリティ分野で活躍したいという意志をもった若者が、合宿形式で実践的な知識を学ぶ。}
  {プログラマーやデザイナーなどから成る複数の参加チームが、与えられたテーマに関するプロトタイプを短期間で作成し、その成果を発表して競い合う。}
  {問題解決や利用者獲得などゲーム的な要素のない分野に、デジタル技術を活用したゲームの要素を取り入れることによって、利用者の参加を動機づける。}
\end{itpmon}
\begin{kaisetsu}{ウ}
\textbf{ハッカソン} = hack + marathon。短期集中でプロトタイプを作り競い合うイベント。
アはディベート、イはセキュリティキャンプ、エはゲーミフィケーションの説明。
\end{kaisetsu}

\begin{itpmon}[\sourcebadge{R6}\enspace\diffbadge{標}]{%
未来のある時点に目標を設定し、そこを起点に現在を振り返り、目標実現のために現在すべきことを考える方法を表す用語として、最も適切なものはどれか。}
\sentakushi
  {PoC（Proof of Concept）}
  {PoV（Proof of Value）}
  {バックキャスティング}
  {フォアキャスティング}
\end{itpmon}
\begin{kaisetsu}{ウ}
\textbf{バックキャスティング}は未来の目標から逆算して現在の行動を決める手法。SDGsなどの長期目標設定で使われる。フォアキャスティングは現状から積み上げる方法で、方向が逆。
\end{kaisetsu}

\begin{itpmon}[\sourcebadge{R6}\enspace\diffbadge{標}]{%
ベンチャーキャピタルに関する記述として、最も適切なものはどれか。}
\sentakushi
  {新しい技術の獲得や、規模の経済性の追求などを目的に、他の企業と共同出資会社を設立する手法}
  {株式売却による利益獲得などを目的に、新しい製品やサービスを武器に市場に参入しようとする企業に対して出資などを行う企業}
  {新サービスや技術革新などの創出を目的に、国や学術機関、他の企業など外部の組織と共創関係を結び、積極的に技術や資源を交換し、自社に取り込む手法}
  {特定された課題の解決を目的に、一定の期間を定めて企業内に立ち上げられ、構成員を関連部門から招集し、目的が達成された時点で解散する組織}
\end{itpmon}
\begin{kaisetsu}{イ}
\textbf{ベンチャーキャピタル（VC）}は成長見込みのあるスタートアップに出資し、株式売却（IPO等）で利益を得る投資会社。
アはジョイントベンチャー、ウはオープンイノベーション、エはタスクフォースの説明。
\end{kaisetsu}

\begin{itpmon}[\sourcebadge{R6}\enspace\diffbadge{強}]{%
企業の戦略立案やマーケティングなどで使用されるフェルミ推定に関する記述として、最も適切なものはどれか。}
\sentakushi
  {正確に算出することが極めて難しい数量に対して、把握している情報と論理的な思考プロセスによって概数を求める手法である。}
  {特定の集団と活動を共にしたり、人々の動きを観察したりすることによって、慣習や嗜好、地域や組織を取り巻く文化を類推する手法である。}
  {入力データと出力データから、その因果関係を統計的に推定する手法である。}
  {有識者のグループに繰り返し同一のアンケート調査とその結果のフィードバックを行うことによって、ある分野の将来予測に関する総意を得る手法である。}
\end{itpmon}
\begin{kaisetsu}{ア}
\textbf{フェルミ推定}は、正確な数値を得ることが困難な量を論理的な分解と概算で推定する手法。
イはエスノグラフィ、ウは回帰分析、エはデルファイ法の説明。
\end{kaisetsu}

% =============================================================
\section{タイムアタック}
% =============================================================

\begin{timeattack}{ストラテジ用語クイックチェック}{5分}{6問中5問}

\begin{itpmon}[\diffbadge{基}]{PPMで「市場成長率が低く、市場占有率が高い」事業を何と呼ぶか。}
\sentakushi{花形}{金のなる木}{問題児}{負け犬}
\end{itpmon}
\begin{kaisetsu}{イ}
成長率低・占有率高 = \textbf{金のなる木}。投資を抑えて利益を回収するフェーズ。
\end{kaisetsu}

\begin{itpmon}[\diffbadge{基}]{SWOT分析で扱うのはどれか。}
\sentakushi
  {市場の成長率と占有率}
  {自社の強み・弱みと、外部の機会・脅威}
  {製品・価格・流通・販促の4要素}
  {顧客・競合・自社の3視点}
\end{itpmon}
\begin{kaisetsu}{イ}
アはPPM、ウは4P、エは3C。SWOTは「内部（強み弱み）×外部（機会脅威）」の4象限。
\end{kaisetsu}

\begin{itpmon}[\diffbadge{基}]{DXの説明として最も適切なものはどれか。}
\sentakushi
  {MVPで仮説検証を繰り返す手法}
  {ITを使って経営やビジネスモデルを変革すること}
  {短期集中でプロトタイプを作り競う}
  {概念や技術の実現可能性を検証すること}
\end{itpmon}
\begin{kaisetsu}{イ}
DX = Digital Transformation。アはリーンスタートアップ、ウはハッカソン、エはPoC。
\end{kaisetsu}

\begin{itpmon}[\diffbadge{基}]{PoCとPoVの正しい関係はどれか。}
\sentakushi
  {PoCは価値の実証、PoVは概念の実証}
  {PoCは概念の実証、PoVは価値の実証}
  {どちらも同じ意味}
  {PoCは市場調査、PoVは技術検証}
\end{itpmon}
\begin{kaisetsu}{イ}
PoC = Proof of \textbf{Concept}（概念実証）。PoV = Proof of \textbf{Value}（価値実証）。PoCが先、PoVが後。
\end{kaisetsu}

\begin{itpmon}[\diffbadge{標}]{フォーラム標準の説明として適切なものはどれか。}
\sentakushi
  {公的機関が策定した規格}
  {市場で事実上の標準となったもの}
  {特定分野の企業が集まって策定した規格}
  {個人が提案して広まったもの}
\end{itpmon}
\begin{kaisetsu}{ウ}
フォーラム標準 = 業界の企業団体が合意して作る規格。アはデジュール標準、イはデファクト標準。
\end{kaisetsu}

\begin{itpmon}[\sourcebadge{R6}\enspace\diffbadge{標}]{%
技術戦略の策定や技術開発の推進といった技術経営に直接の責任をもつ役職はどれか。}
\sentakushi{CEO}{CFO}{COO}{CTO}
\end{itpmon}
\begin{kaisetsu}{エ}
\textbf{CTO}（Chief Technology Officer）= 最高技術責任者。技術経営に直接の責任を持つ。
\end{kaisetsu}

\end{timeattack}

% =============================================================
\section{よくあるミスと到達確認}
% =============================================================

\begin{yokumiss}[この章でよくあるミス]
\begin{enumerate}[leftmargin=1.6em, itemsep=5pt]
  \item \textbf{SWOTとVRIOを混同する}——SWOTは4象限マトリクス、VRIOは強みの深堀り。目的が違う。
  \item \textbf{PoCとMVPを同一視する}——PoCは「技術の実現可能性確認」、MVPは「最小限の製品で市場仮説を検証」。
  \item \textbf{DXとリーンスタートアップを混同する}——DXは概念、リーンスタートアップは手法。
  \item \textbf{CIOとCTOを逆にする}——CIOは「情報システム戦略」、CTOは「技術開発」。
\end{enumerate}
\end{yokumiss}

\begin{tatsaku}
  \item 経営理念→方針→戦略→戦術の4段階を説明できる
  \item CxOの各役職の守備範囲を即答できる
  \item SWOT/VRIO/3C/5F/4P/PPMの「何を見る道具か」を言える
  \item PPMの4象限を図なしで再現できる
  \item PoC/PoV/MVP/リーンスタートアップの違いを説明できる
  \item バックキャスティングとフォアキャスティングの違いを言える
  \item ハッカソン/ゲーミフィケーション/セキュリティキャンプを区別できる
  \item ベンチャーキャピタルとジョイントベンチャーの違いを説明できる
\end{tatsaku}

\nextchapter{%
会計・財務・法務を学びます。PLの5段階利益、損益分岐点、知的財産権、契約形態——
ストラテジ系の後半を一気に攻略します。
}

% =============================================================
% 02_accounting.tex — 会計・財務と法務
% =============================================================
\chapter{会計・財務と法務}

\chaptermeta{%
  \item PLの5段階利益を読み取れる
  \item 損益分岐点（BEP）を計算できる
  \item 知的財産権の種類と保護期間を整理できる
  \item 請負・準委任・派遣の違いを即答できる
  \item 情報倫理・コンプライアンスの基本を押さえる
}{75分}{2}{第1章}

ストラテジ系の後半は「お金」と「法律」です。
PLの5段階利益、損益分岐点の計算、知的財産権、契約形態——
\textbf{数字の出る問題}と\textbf{法律の識別問題}を確実に得点源にしましょう。

% =============================================================
\section{財務三表の全体像}
% =============================================================

\begin{teigi}[財務三表とは]
企業の経営状態を把握するための3つの書類。
\begin{itemize}
  \item \termtag{損益計算書（PL）}：一定期間の「儲け」を示す（フロー）
  \item \termtag{貸借対照表（BS）}：ある時点の「財産」を示す（ストック）
  \item \termtag{キャッシュフロー計算書（CF）}：資金の流れを示す
\end{itemize}
Iパスでは特にPLの出題頻度が高い。BSは「資産＝負債＋純資産」を押さえればOK。
\end{teigi}

\begin{pitfall}[BSの左右を逆にしがち]
貸借対照表は\textbf{左側（借方）＝資産}、\textbf{右側（貸方）＝負債＋純資産}。
「左に持っているもの、右にその調達元」と覚える。
\end{pitfall}

% =============================================================
\section{PLの5段階利益}
% =============================================================

\begin{teigi}[PLの5段階利益]
PLは売上高から順に費用を引いて5つの利益を計算する。

\begin{enumerate}
  \item \termtag{売上総利益}＝ 売上高 $-$ 売上原価（粗利）
  \item \termtag{営業利益}＝ 売上総利益 $-$ 販管費（本業の儲け）
  \item \termtag{経常利益}＝ 営業利益 $+$ 営業外収益 $-$ 営業外費用
  \item \termtag{税引前当期純利益}＝ 経常利益 $+$ 特別利益 $-$ 特別損失
  \item \termtag{当期純利益}＝ 税引前当期純利益 $-$ 法人税等
\end{enumerate}
\end{teigi}

\begin{hikaku}[PLの5段階を一覧比較]
\comptable{利益 & 計算式 & 何がわかるか}{
売上総利益 & 売上高 $-$ 売上原価 & 商品力 \\
営業利益 & 売上総利益 $-$ 販管費 & 本業の収益力 \\
経常利益 & 営業利益 $\pm$ 営業外損益 & 通常の経営力 \\
税引前当期純利益 & 経常利益 $\pm$ 特別損益 & 臨時含む全体像 \\
当期純利益 & 税引前 $-$ 法人税等 & 最終的な儲け
}
\end{hikaku}

% --- 例題 ---
\begin{reidai}{PLの計算（当期純利益から販管費を逆算）}
ある企業のPLが以下のとおりであるとき、販売費及び一般管理費は何百万円か。

\begin{center}
\small
\begin{tabular}{lr}
\toprule
売上高 & 5,000百万円 \\
売上原価 & 3,000百万円 \\
営業外収益 & 100百万円 \\
営業外費用 & 200百万円 \\
特別利益 & 50百万円 \\
特別損失 & 50百万円 \\
法人税等 & 100百万円 \\
当期純利益 & 800百万円 \\
\bottomrule
\end{tabular}
\end{center}

\sentakushi{800}{900}{1,000}{1,100}
\end{reidai}

\begin{shishin}
当期純利益から逆算する。税引前当期純利益 $=$ 当期純利益 $+$ 法人税等 $= 800 + 100 = 900$。
経常利益 $=$ 税引前 $-$ 特別利益 $+$ 特別損失 $= 900 - 50 + 50 = 900$。
営業利益 $=$ 経常利益 $-$ 営業外収益 $+$ 営業外費用 $= 900 - 100 + 200 = 1{,}000$。
売上総利益 $=$ 売上高 $-$ 売上原価 $= 5{,}000 - 3{,}000 = 2{,}000$。
販管費 $=$ 売上総利益 $-$ 営業利益 $= 2{,}000 - 1{,}000 = 1{,}000$。
\end{shishin}

\begin{kaitou}
\textbf{正解：ウ（1,000百万円）}

5段階利益を下から逆算する典型パターン。
当期純利益→税引前→経常→営業の順に巻き戻し、最後に売上総利益との差で販管費を求める。
\end{kaitou}

\begin{minuki}[見抜き方]
\begin{itemize}
  \item PLの問題は「逆算」がほとんど——与えられた利益から上方向に戻して未知を求める
  \item 特別損益がゼロ（相殺）なら経常利益＝税引前当期純利益と気づけるとスピードアップ
  \item 営業外損益の符号を間違えやすい——「収益はプラス、費用はマイナス」を徹底
\end{itemize}
\end{minuki}

% =============================================================
\section{損益分岐点}
% =============================================================

\begin{teigi}[損益分岐点（BEP）]
利益がちょうどゼロになる売上高。
\[
  \text{BEP} = \frac{\text{固定費}}{1 - \text{変動費率}}
  = \frac{\text{固定費}}{\text{限界利益率}}
\]
\begin{itemize}
  \item \termtag{固定費}：売上に関係なく一定（家賃、人件費など）
  \item \termtag{変動費}：売上に比例して増減（材料費など）
  \item \termtag{変動費率}＝ 変動費 ÷ 売上高
  \item \termtag{限界利益率}＝ 1 $-$ 変動費率
\end{itemize}
\end{teigi}

\begin{pitfall}[固定費と変動費の分類ミス]
\begin{itemize}
  \item 人件費は\textbf{固定費}（パート・アルバイトの時給制でも基本は固定費扱い）
  \item 材料費・仕入原価は\textbf{変動費}
  \item 問題文の「売上に比例する」かどうかで判断
\end{itemize}
\end{pitfall}

% =============================================================
\section{知的財産権}
% =============================================================

\begin{hikaku}[知的財産権の比較]
\comptable{権利 & 保護対象 & 保護期間}{
特許権 & 発明（技術的アイデア） & 出願から20年 \\
実用新案権 & 考案（物品の形状等） & 出願から10年 \\
意匠権 & デザイン & 出願から25年 \\
商標権 & ブランド名・ロゴ & 登録から10年（更新可） \\
著作権 & 創作的な表現 & 著作者の死後70年
}
\end{hikaku}

\begin{teigi}[著作権の特徴]
\begin{itemize}
  \item \termtag{無方式主義}：登録不要で、創作した時点で自動的に発生
  \item プログラムも著作権で保護される（ただしアルゴリズムは保護されない）
  \item アイデアそのものは著作権では保護されない（表現が対象）
\end{itemize}
\end{teigi}

\begin{pitfall}[特許権と著作権の混同]
\begin{itemize}
  \item \textbf{特許権}：出願・審査が必要（方式主義）、技術的アイデアを保護
  \item \textbf{著作権}：登録不要（無方式主義）、表現を保護
  \item 「プログラムの著作権」と「ソフトウェア特許」は別物
\end{itemize}
\end{pitfall}

% =============================================================
\section{契約形態（請負・準委任・派遣）}
% =============================================================

\begin{hikaku}[3つの契約形態の比較]
\comptable{項目 & 請負 & 準委任 & 派遣}{
完成責任 & あり（成果物を納品） & なし（善管注意義務） & なし \\
指揮命令 & 受注者 & 受注者 & 派遣先 \\
瑕疵担保 & あり & なし & --- \\
報酬の対象 & 成果物 & 作業時間・工数 & 労働時間
}
\end{hikaku}

\begin{teigi}[偽装請負とは]
\termtag{偽装請負}：形式上は請負契約だが、実態は発注者が労働者に直接指揮命令している状態。
\textbf{違法}であり、労働者派遣法に違反する。
\begin{itemize}
  \item 請負なのに発注者が作業の進め方を細かく指示 → 偽装請負
  \item 指揮命令権が発注者にあるなら、正しくは「派遣契約」を結ぶべき
\end{itemize}
\end{teigi}

% =============================================================
\section{情報倫理とコンプライアンス}
% =============================================================

\begin{teigi}[主要な法律]
\begin{itemize}
  \item \termtag{個人情報保護法}：個人情報の取扱いルール（利用目的の特定・公表）
  \item \termtag{不正アクセス禁止法}：他人のID/パスワードの不正使用を禁止
  \item \termtag{特定電子メール法}：広告宣伝メールのオプトイン規制
  \item \termtag{不正競争防止法}：営業秘密の不正取得・ドメイン名の不正使用を禁止
\end{itemize}
\end{teigi}

% =============================================================
\section{過去問ドリル}
% =============================================================

\shoumatsuhead{過去問ドリル：会計・法務}

\begin{itpmon}[\sourcebadge{R7}\enspace\diffbadge{標}]{%
A社がシステム開発をB社に委託している。A社の社員がB社の作業者に対して、日々の作業の割振りや進捗管理を直接行っている場合、この状態を何というか。}
\sentakushi
  {SES契約}
  {偽装請負}
  {派遣契約}
  {準委任契約}
\end{itpmon}
\begin{kaisetsu}{イ}
請負や準委任では指揮命令権は受注者側にある。発注者が直接作業指示をしている実態があれば、形式に関わらず\textbf{偽装請負}にあたる。ウの派遣契約であれば合法だが、請負契約のまま指揮命令していることが問題。
\end{kaisetsu}

\begin{itpmon}[\sourcebadge{R7}\enspace\diffbadge{標}]{%
特定電子メール法で規制される行為として、最も適切なものはどれか。}
\sentakushi
  {企業が自社の社員全員に業務連絡のメールを送信した。}
  {受信者の同意を得ずに、広告宣伝の電子メールを送信した。}
  {知人に転職の挨拶メールを送信した。}
  {取引先に見積書をメールで送信した。}
\end{itpmon}
\begin{kaisetsu}{イ}
\textbf{特定電子メール法}は「広告宣伝メール」に対して\textbf{オプトイン規制}（事前同意なしの送信を禁止）を定めている。業務連絡や個人的な連絡は規制対象外。
\end{kaisetsu}

\begin{itpmon}[\sourcebadge{R6}\enspace\diffbadge{標}]{%
著作権に関する記述として、最も適切なものはどれか。}
\sentakushi
  {著作権を取得するには、特許庁に登録する必要がある。}
  {著作権は、著作物を創作した時点で自動的に発生する。}
  {アルゴリズムは著作権で保護される。}
  {著作権の保護期間は、著作者の死後100年である。}
\end{itpmon}
\begin{kaisetsu}{イ}
著作権は\textbf{無方式主義}——登録不要で創作時に自動発生する。アは特許権の手続き。ウについてアルゴリズム（手順）は著作権の対象外。エの保護期間は死後\textbf{70年}。
\end{kaisetsu}

% =============================================================
\section{タイムアタック}
% =============================================================

\begin{timeattack}{会計・法務クイックチェック}{5分}{4問中3問}

\begin{itpmon}[\diffbadge{基}]{PLにおいて「売上高 $-$ 売上原価」で求められる利益はどれか。}
\sentakushi{営業利益}{経常利益}{売上総利益}{当期純利益}
\end{itpmon}
\begin{kaisetsu}{ウ}
売上高 $-$ 売上原価 $=$ \textbf{売上総利益}（粗利）。営業利益はさらに販管費を引いたもの。
\end{kaisetsu}

\begin{itpmon}[\diffbadge{標}]{固定費が600万円、変動費率が40\%のとき、損益分岐点売上高は何万円か。}
\sentakushi{800}{1,000}{1,200}{1,500}
\end{itpmon}
\begin{kaisetsu}{イ}
BEP $= 600 \div (1 - 0.4) = 600 \div 0.6 = 1{,}000$万円。
\end{kaisetsu}

\begin{itpmon}[\diffbadge{基}]{出願から20年間保護される知的財産権はどれか。}
\sentakushi{意匠権}{実用新案権}{商標権}{特許権}
\end{itpmon}
\begin{kaisetsu}{エ}
\textbf{特許権}は出願から20年。意匠権は25年、実用新案権は10年、商標権は登録から10年（更新可）。
\end{kaisetsu}

\begin{itpmon}[\diffbadge{標}]{請負契約と準委任契約の最も大きな違いはどれか。}
\sentakushi
  {契約の相手方が法人か個人かの違い}
  {成果物の完成責任があるかないかの違い}
  {作業場所がどこかの違い}
  {報酬が月額か日額かの違い}
\end{itpmon}
\begin{kaisetsu}{イ}
請負は\textbf{成果物の完成責任}がある。準委任は善管注意義務はあるが完成責任はない。
\end{kaisetsu}

\end{timeattack}

% =============================================================
\section{よくあるミスと到達確認}
% =============================================================

\begin{yokumiss}[この章でよくあるミス]
\begin{enumerate}[leftmargin=1.6em, itemsep=5pt]
  \item \textbf{営業利益と経常利益を混同する}——営業利益は本業の儲け、経常利益は営業外損益を含む。
  \item \textbf{損益分岐点の公式で変動費率と限界利益率を逆にする}——分母は$1-$変動費率（＝限界利益率）。
  \item \textbf{特許権と著作権の取得方法を混同する}——特許は出願・審査が必要、著作権は自動発生。
  \item \textbf{請負と準委任の指揮命令権を間違える}——どちらも指揮命令は受注者。派遣だけが派遣先。
  \item \textbf{意匠権の保護期間を忘れる}——2020年改正で出願から25年に延長された。
\end{enumerate}
\end{yokumiss}

\begin{tatsaku}
  \item PLの5段階利益を順番に言える
  \item 損益分岐点の公式を使って計算できる
  \item 知的財産権5種の保護対象と期間を言える
  \item 著作権の無方式主義を説明できる
  \item 請負・準委任・派遣の3点比較表を再現できる
  \item 偽装請負の判断ポイントを説明できる
  \item 特定電子メール法のオプトイン規制を理解している
  \item BSの左右（資産＝負債＋純資産）を説明できる
\end{tatsaku}

\nextchapter{%
マネジメント系に入ります。システム開発の5工程、ウォーターフォールとアジャイルの比較、
プロジェクトマネジメントのQCD——「つくる側」の視点を学びます。
}


% ── PART II: マネジメント系 ──
\partpage{II}{マネジメント系}
% =============================================================
% 03_development.tex — システム開発とプロジェクトマネジメント
% =============================================================
\chapter{システム開発とプロジェクトマネジメント}

\chaptermeta{%
  \item 開発工程の5段階を順序どおり説明できる
  \item ウォーターフォールとアジャイルの違いを即答できる
  \item QCD（品質・コスト・納期）のトレードオフを理解できる
  \item WBS/ガントチャートを使い分けられる
  \item テスト工程のV字モデルを説明できる
}{75分}{2}{第1--2章}

マネジメント系の前半は「つくる側」の話です。
どんな順番で開発し、どう管理するか——
\textbf{開発モデルの識別}と\textbf{PMの基本用語}を押さえれば得点源になります。

% =============================================================
\section{開発工程の5段階}
% =============================================================

\begin{teigi}[共通フレーム（SLCP）の5工程]
\begin{enumerate}
  \item \termtag{要件定義}：「何を作るか」を決める（ユーザ主体）
  \item \termtag{外部設計（基本設計）}：画面・帳票・インターフェースを決める
  \item \termtag{内部設計（詳細設計）}：プログラムの内部構造を決める
  \item \termtag{プログラミング}：実際にコードを書く
  \item \termtag{テスト}：単体→結合→システム→運用テストの順に実施
\end{enumerate}
\end{teigi}

\begin{pitfall}[要件定義と外部設計を混同しがち]
\begin{itemize}
  \item \textbf{要件定義}は「何をしたいか」を洗い出す（What）
  \item \textbf{外部設計}は「どう見せるか」を決める（How の外側）
  \item 要件定義はユーザ主体、外部設計は開発者主体
\end{itemize}
\end{pitfall}

\subsection{テスト工程とV字モデル}

\begin{hikaku}[V字モデル：設計とテストの対応]
\comptable{設計工程 & 対応するテスト & 確認内容}{
要件定義 & 運用テスト（承認テスト） & ユーザの要件を満たすか \\
外部設計 & システムテスト & システム全体が正しく動くか \\
内部設計 & 結合テスト & モジュール間の連携 \\
プログラミング & 単体テスト & 個々のモジュールの正しさ
}
\end{hikaku}

% =============================================================
\section{開発モデルの比較}
% =============================================================

\begin{hikaku}[代表的な開発モデル]
\comptable{モデル & 特徴 & 向いている場面}{
ウォーターフォール & 工程を順に進める・後戻りしにくい & 要件が明確な大規模開発 \\
アジャイル & 短い反復で少しずつ完成 & 要件変更が多いプロジェクト \\
プロトタイピング & 試作品で要件を確認 & ユーザのイメージが曖昧 \\
スパイラル & 設計→試作→評価を繰り返す & リスクの高い大規模開発 \\
DevOps & 開発と運用の一体化 & 継続的デリバリーが必要
}
\end{hikaku}

\begin{teigi}[アジャイルのキーワード]
\begin{itemize}
  \item \termtag{イテレーション}：1〜4週間の短い開発サイクル
  \item \termtag{スクラム}：代表的なアジャイルフレームワーク
  \item \termtag{スプリント}：スクラムにおけるイテレーション
  \item \termtag{プロダクトオーナー}：要件の優先順位を決める人
  \item \termtag{デイリースクラム}：毎日の短い進捗共有ミーティング
\end{itemize}
\end{teigi}

% --- 例題 ---
\begin{reidai}{開発モデルの識別}
システム開発において、短い期間で設計・実装・テストを繰り返し、
少しずつ機能を追加していく開発手法はどれか。

\sentakushi{ウォーターフォールモデル}{アジャイル開発}{共通フレーム}{PMBOK}
\end{reidai}

\begin{shishin}
「短い期間」「繰り返し」「少しずつ機能を追加」がキーワード。
ウォーターフォールは「順に進める」、共通フレームは「開発プロセスの規格」、
PMBOKは「PM知識体系」なので該当しない。
\end{shishin}

\begin{kaitou}
\textbf{正解：イ（アジャイル開発）}

短いイテレーションで反復的に開発する手法がアジャイル開発。
ウォーターフォールは工程を順に進め、後戻りしにくい点が対照的。
\end{kaitou}

\begin{minuki}[見抜き方]
\begin{itemize}
  \item 「繰り返し」「反復」「短期間」→ アジャイル
  \item 「順番に」「後戻りしない」「前工程完了後に次へ」→ ウォーターフォール
  \item 「試作品を見せて確認」→ プロトタイピング
  \item 「開発と運用の一体化」「CI/CD」→ DevOps
\end{itemize}
\end{minuki}

% =============================================================
\section{プロジェクトマネジメントとQCD}
% =============================================================

\begin{teigi}[QCD：プロジェクトの三大制約]
\begin{itemize}
  \item \termtag{Q（Quality）}：品質——成果物の正しさ・完成度
  \item \termtag{C（Cost）}：コスト——予算・工数
  \item \termtag{D（Delivery）}：納期——スケジュール
\end{itemize}
3つはトレードオフの関係にあり、1つを優先すると他が犠牲になりやすい。
\end{teigi}

\begin{pitfall}[「ファンクションポイント法」と「COCOMO」の混同]
\begin{itemize}
  \item \textbf{ファンクションポイント法}：機能の数と複雑さで工数を見積もる
  \item \textbf{COCOMO}：プログラムの行数で工数を見積もる
  \item 「機能」と「行数」のどちらに着目するかが違う
\end{itemize}
\end{pitfall}

% =============================================================
\section{WBS・ガントチャート}
% =============================================================

\begin{teigi}[WBSとガントチャート]
\begin{itemize}
  \item \termtag{WBS}（Work Breakdown Structure）：作業を階層的に分解した構造図。プロジェクトの全作業を洗い出す。
  \item \termtag{ガントチャート}：横軸に時間、縦軸に作業項目をとった棒グラフ。進捗管理に使う。
  \item \termtag{アローダイアグラム}：作業の依存関係を矢印で表す。\textbf{クリティカルパス}（最長経路）を見つけて全体の最短所要日数を求める。
\end{itemize}
\end{teigi}

\begin{hikaku}[PM手法の使い分け]
\comptable{ツール & 目的 & 特徴}{
WBS & 作業の洗い出し & ツリー構造で漏れなく分解 \\
ガントチャート & 進捗の見える化 & 横棒で期間を表示 \\
アローダイアグラム & 依存関係と最短日数 & クリティカルパスの特定 \\
PERT & 所要時間の見積もり & 楽観/最頻/悲観の3点見積もり
}
\end{hikaku}

% =============================================================
\section{過去問ドリル}
% =============================================================

\shoumatsuhead{過去問ドリル：開発とPM}

\begin{itpmon}[\sourcebadge{R6}\enspace\diffbadge{標}]{%
業務プロセスを可視化する手法のうち、業務の流れを「イベント」「アクティビティ」「ゲートウェイ」などの記号で表記するものはどれか。}
\sentakushi
  {BPMN（Business Process Model and Notation）}
  {DFD（Data Flow Diagram）}
  {E-R図}
  {UML}
\end{itpmon}
\begin{kaisetsu}{ア}
\textbf{BPMN}は業務プロセスをイベント・アクティビティ・ゲートウェイの記号で可視化する標準表記法。DFDはデータの流れ、E-R図はデータの関係、UMLはオブジェクト指向の設計図。
\end{kaisetsu}

\begin{itpmon}[\sourcebadge{original}\enspace\diffbadge{標}]{%
アジャイル開発の特徴として、最も適切なものはどれか。}
\sentakushi
  {全工程の計画を詳細に立ててから開発に着手する。}
  {短い期間での反復を繰り返し、段階的にシステムを完成させる。}
  {利用者にプロトタイプを見せて要件を確定する。}
  {プログラムのステップ数から開発コストを見積もる。}
\end{itpmon}
\begin{kaisetsu}{イ}
アジャイルは短いイテレーションで反復的に開発する手法。アはウォーターフォール、ウはプロトタイピング、エはCOCOMOの説明。
\end{kaisetsu}

\begin{itpmon}[\sourcebadge{original}\enspace\diffbadge{標}]{%
WBS（Work Breakdown Structure）の説明として、最も適切なものはどれか。}
\sentakushi
  {プロジェクトの作業を階層的に分解して構造化したもの}
  {横軸に時間、縦軸に作業をとって進捗を管理する図表}
  {作業の依存関係を矢印で表し、クリティカルパスを求める手法}
  {品質・コスト・納期の三大制約を管理するフレームワーク}
\end{itpmon}
\begin{kaisetsu}{ア}
\textbf{WBS}はプロジェクトの全作業を階層的に分解した構造。イはガントチャート、ウはアローダイアグラム（PERT図）、エはQCDの説明。
\end{kaisetsu}

% =============================================================
\section{タイムアタック}
% =============================================================

\begin{timeattack}{開発・PMクイックチェック}{5分}{4問中3問}

\begin{itpmon}[\diffbadge{基}]{システム開発の工程で、最初に行うのはどれか。}
\sentakushi{外部設計}{テスト}{プログラミング}{要件定義}
\end{itpmon}
\begin{kaisetsu}{エ}
開発は\textbf{要件定義}→外部設計→内部設計→プログラミング→テストの順。
\end{kaisetsu}

\begin{itpmon}[\diffbadge{基}]{V字モデルで、要件定義に対応するテストはどれか。}
\sentakushi{単体テスト}{結合テスト}{システムテスト}{運用テスト}
\end{itpmon}
\begin{kaisetsu}{エ}
要件定義 ↔ \textbf{運用テスト}。外部設計 ↔ システムテスト、内部設計 ↔ 結合テスト、プログラミング ↔ 単体テスト。
\end{kaisetsu}

\begin{itpmon}[\diffbadge{標}]{スクラムにおける「スプリント」の説明として適切なものはどれか。}
\sentakushi
  {プロジェクト全体の工程を順に進める方式}
  {1〜4週間の短い反復開発の単位}
  {開発と運用を一体化する取組み}
  {試作品を使って要件を確認する手法}
\end{itpmon}
\begin{kaisetsu}{イ}
\textbf{スプリント}はスクラムにおける1〜4週間のイテレーション。アはウォーターフォール、ウはDevOps、エはプロトタイピング。
\end{kaisetsu}

\begin{itpmon}[\diffbadge{標}]{プロジェクトの全体所要日数を求めるために使う「クリティカルパス」とは何か。}
\sentakushi
  {コストが最も大きい作業の経路}
  {品質リスクが最も高い作業の経路}
  {所要時間が最も長い作業の経路}
  {担当者が最も多い作業の経路}
\end{itpmon}
\begin{kaisetsu}{ウ}
\textbf{クリティカルパス}はアローダイアグラムにおける最長経路であり、プロジェクト全体の最短所要日数を決定する。
\end{kaisetsu}

\end{timeattack}

% =============================================================
\section{よくあるミスと到達確認}
% =============================================================

\begin{yokumiss}[この章でよくあるミス]
\begin{enumerate}[leftmargin=1.6em, itemsep=5pt]
  \item \textbf{V字モデルの対応関係を間違える}——「一番上同士、一番下同士」で対応する。
  \item \textbf{アジャイルを「計画しない開発」と誤解する}——計画はする。ただし短いサイクルで柔軟に変更する。
  \item \textbf{WBSとガントチャートを混同する}——WBSは「分解」、ガントチャートは「日程管理」。
  \item \textbf{ファンクションポイント法とCOCOMOの違いを忘れる}——「機能数」か「行数」か。
  \item \textbf{クリティカルパスを「最短経路」と勘違いする}——最長経路が正解。最長経路＝全体の最短所要日数。
\end{enumerate}
\end{yokumiss}

\begin{tatsaku}
  \item 開発工程の5段階を順番に言える
  \item V字モデルの対応関係を4組すべて言える
  \item ウォーターフォールとアジャイルの違いを3点以上挙げられる
  \item アジャイル関連用語（スクラム/スプリント/PO）を説明できる
  \item QCDの三大制約とトレードオフを説明できる
  \item WBS/ガントチャート/アローダイアグラムの目的を区別できる
  \item クリティカルパスの意味と求め方を説明できる
  \item BPMNとDFDの違いを説明できる
\end{tatsaku}

\nextchapter{%
サービスマネジメントと監査に進みます。SLA/SLM、ITIL、稼働率の計算、
システム監査の独立性——「運用する側」の視点を学びます。
}

% =============================================================
% 04_service.tex — サービスマネジメントと監査
% =============================================================
\chapter{サービスマネジメントと監査}

\chaptermeta{%
  \item SLAとSLMの違いを即答できる
  \item ITILのサービス管理プロセスを整理できる
  \item MTBF/MTTRから稼働率を計算できる
  \item システム監査の独立性を説明できる
  \item 内部統制の目的と構成要素を押さえる
}{60分}{2}{第3章}

マネジメント系の後半は「運用する側」の話です。
作ったシステムをどう維持し、どう監査するか——
\textbf{計算問題}が出る稼働率と、\textbf{概念の識別}が出るSLA/監査を中心に学びます。

% =============================================================
\section{SLA/SLM}
% =============================================================

\begin{teigi}[SLAとSLM]
\begin{itemize}
  \item \termtag{SLA}（Service Level Agreement）：サービス提供者と利用者の間で合意した\textbf{サービス品質の取り決め}。稼働率99.9\%、応答時間3秒以内、といった具体的な数値で定義する。
  \item \termtag{SLM}（Service Level Management）：SLAで定めたサービスレベルを\textbf{継続的に管理・改善するプロセス}。
\end{itemize}
SLAは「約束（文書）」、SLMは「約束を守る活動（プロセス）」。
\end{teigi}

\begin{pitfall}[SLAとSLMを逆にしがち]
\begin{itemize}
  \item SLA → \textbf{A}greement = 合意 = 文書
  \item SLM → \textbf{M}anagement = 管理 = プロセス
  \item 「Aは文書、Mは活動」と頭文字で区別
\end{itemize}
\end{pitfall}

% =============================================================
\section{サービスデスクとITIL}
% =============================================================

\begin{teigi}[ITILとサービスデスク]
\begin{itemize}
  \item \termtag{ITIL}：ITサービスマネジメントのベストプラクティス集
  \item \termtag{サービスデスク}：利用者からの問い合わせを受ける単一の窓口（SPOC）
  \item \termtag{インシデント管理}：サービスの中断や品質低下に迅速に対応し、正常状態に戻す
  \item \termtag{問題管理}：インシデントの根本原因を特定し、再発を防止する
  \item \termtag{変更管理}：ITサービスへの変更を管理し、リスクを最小化する
\end{itemize}
\end{teigi}

\begin{hikaku}[インシデント管理と問題管理]
\comptable{項目 & インシデント管理 & 問題管理}{
目的 & サービスの早期復旧 & 根本原因の除去 \\
時間軸 & 短期（今すぐ対処） & 中長期（再発防止） \\
例え & 応急処置 & 根治治療
}
\end{hikaku}

\begin{pitfall}[インシデント管理で「原因究明」を選びがち]
インシデント管理の目的は\textbf{迅速な復旧}。原因究明は問題管理の役割。
「まず直す（インシデント）→ 次に原因を探る（問題）」の順。
\end{pitfall}

% =============================================================
\section{稼働率の計算}
% =============================================================

\begin{teigi}[MTBF/MTTR/稼働率]
\begin{itemize}
  \item \termtag{MTBF}（Mean Time Between Failures）：平均故障間隔（動いている時間）
  \item \termtag{MTTR}（Mean Time To Repair）：平均修理時間（止まっている時間）
  \item \termtag{稼働率} $= \dfrac{\text{MTBF}}{\text{MTBF} + \text{MTTR}}$
\end{itemize}
\end{teigi}

% --- 例題 ---
\begin{reidai}{稼働率の計算}
あるサーバのMTBFが900時間、MTTRが100時間であるとき、稼働率はいくらか。

\sentakushi{0.1}{0.8}{0.9}{0.99}
\end{reidai}

\begin{shishin}
公式に当てはめる。稼働率 $= \text{MTBF} \div (\text{MTBF} + \text{MTTR})$。
$= 900 \div (900 + 100) = 900 \div 1{,}000 = 0.9$。
\end{shishin}

\begin{kaitou}
\textbf{正解：ウ（0.9）}

稼働率 $= 900 \div (900 + 100) = 0.9$（90\%）。
MTBFとMTTRを足すと全体の時間（MTBF + MTTR）になり、
そのうちMTBFの割合が稼働率。
\end{kaitou}

\begin{minuki}[見抜き方]
\begin{itemize}
  \item 稼働率の公式は「動いている時間 ÷ 全体の時間」と覚える
  \item 直列システムの稼働率 $=$ 各稼働率の積
  \item 並列システムの稼働率 $= 1 - (1 - A_1)(1 - A_2)$
  \item 直列は「どちらも動く」、並列は「どちらかが動けばOK」
\end{itemize}
\end{minuki}

\subsection{直列・並列システムの稼働率}

\begin{hikaku}[直列と並列の稼働率]
\comptable{構成 & 計算式 & 特徴}{
直列 & $A_1 \times A_2$ & 全部動かないとダメ → 稼働率が下がる \\
並列（冗長化） & $1 - (1-A_1)(1-A_2)$ & 片方動けばOK → 稼働率が上がる
}
\end{hikaku}

% =============================================================
\section{システム監査}
% =============================================================

\begin{teigi}[システム監査の基本]
\begin{itemize}
  \item \termtag{システム監査}：情報システムの信頼性・安全性・効率性を独立した立場で検証・評価すること
  \item \termtag{監査人の独立性}：被監査部門から独立していること（利害関係がないこと）
  \item \termtag{監査の流れ}：計画 → 実施 → 報告 → フォローアップ
\end{itemize}
\end{teigi}

\begin{pitfall}[「監査人が改善を実施する」は誤り]
監査人は\textbf{助言・勧告}を行うが、改善の\textbf{実施は被監査部門}が行う。
監査人が自ら改善すると独立性が失われる。
\end{pitfall}

% =============================================================
\section{内部統制}
% =============================================================

\begin{teigi}[内部統制の4つの目的と6つの構成要素]
\textbf{4つの目的}：
\begin{enumerate}
  \item 業務の有効性・効率性
  \item 財務報告の信頼性
  \item 法令遵守（コンプライアンス）
  \item 資産の保全
\end{enumerate}
\textbf{6つの構成要素}：統制環境、リスクの評価と対応、統制活動、情報と伝達、モニタリング、ITへの対応。
\end{teigi}

\begin{teigi}[職務分掌（SoD）]
\termtag{職務分掌}（Segregation of Duties）：不正やミスを防ぐために、一つの業務を複数の担当者に分離すること。
\begin{itemize}
  \item 例：発注者と検収者を分ける、プログラマとテスタを分ける
  \item 「一人で完結できない」仕組みが不正を抑止する
\end{itemize}
\end{teigi}

% =============================================================
\section{過去問ドリル}
% =============================================================

\shoumatsuhead{過去問ドリル：サービスと監査}

\begin{itpmon}[\sourcebadge{original}\enspace\diffbadge{標}]{%
SLAの説明として、最も適切なものはどれか。}
\sentakushi
  {ITサービスマネジメントのベストプラクティスをまとめた文書}
  {サービス提供者と利用者の間で合意したサービス品質の取り決め}
  {サービスの品質を継続的に管理・改善するプロセス}
  {利用者からの問い合わせを受け付ける単一の窓口}
\end{itpmon}
\begin{kaisetsu}{イ}
\textbf{SLA}はサービスレベルの合意文書。アはITIL、ウはSLM、エはサービスデスク（SPOC）の説明。
\end{kaisetsu}

\begin{itpmon}[\sourcebadge{original}\enspace\diffbadge{標}]{%
DevOpsの説明として、最も適切なものはどれか。}
\sentakushi
  {開発工程を上流から下流へ順番に進め、後戻りしない手法}
  {開発チームと運用チームが連携し、継続的にサービスを提供する取組み}
  {プロジェクトの作業を階層的に分解して管理する手法}
  {短いスプリントで反復開発を行うフレームワーク}
\end{itpmon}
\begin{kaisetsu}{イ}
\textbf{DevOps} = Development + Operations。開発と運用が協力してCI/CDを実現する。アはウォーターフォール、ウはWBS、エはスクラム。
\end{kaisetsu}

\begin{itpmon}[\sourcebadge{original}\enspace\diffbadge{強}]{%
システム監査において、監査人が守るべき原則として最も適切なものはどれか。}
\sentakushi
  {監査対象のシステムの開発に自ら参加し、改善提案を行う。}
  {被監査部門の業務改善を監査人が直接実施する。}
  {被監査部門から独立した立場で客観的に評価を行う。}
  {監査結果を被監査部門の上長にのみ報告し、経営者には報告しない。}
\end{itpmon}
\begin{kaisetsu}{ウ}
システム監査の基本原則は\textbf{独立性}。監査人は被監査部門から独立した立場で評価する。ア・イは独立性に反し、エは経営者への報告が必要。
\end{kaisetsu}

% =============================================================
\section{よくあるミスと到達確認}
% =============================================================

\begin{yokumiss}[この章でよくあるミス]
\begin{enumerate}[leftmargin=1.6em, itemsep=5pt]
  \item \textbf{SLAとSLMを逆に覚える}——AはAgreement（合意文書）、MはManagement（管理プロセス）。
  \item \textbf{インシデント管理で「根本原因を除去」を選ぶ}——それは問題管理の目的。
  \item \textbf{稼働率の公式でMTTRを分子にする}——分子はMTBF（動いている時間）。
  \item \textbf{並列の稼働率を「足し算」で計算する}——$1-(1-A_1)(1-A_2)$が正しい。
  \item \textbf{監査人が改善を「実施する」と答える}——監査人は助言・勧告まで。実施は被監査部門。
\end{enumerate}
\end{yokumiss}

\begin{tatsaku}
  \item SLAとSLMの違いを即答できる
  \item インシデント管理と問題管理の違いを説明できる
  \item MTBF/MTTRの意味と稼働率の計算ができる
  \item 直列・並列システムの稼働率を計算できる
  \item ITILの主要プロセスを3つ以上言える
  \item システム監査の独立性を説明できる
  \item 内部統制の4つの目的を言える
  \item 職務分掌の意味と目的を説明できる
\end{tatsaku}

\nextchapter{%
テクノロジ系に入ります。2進数変換、CPU/メモリの階層構造、RAID——
コンピュータの「中身」を理解してITの土台を固めます。
}


% ── PART III: テクノロジ系 ──
\partpage{III}{テクノロジ系}
% =============================================================
% 05_foundation.tex — 基礎理論・ハードウェア・ソフトウェア
% =============================================================
\chapter{基礎理論・ハードウェア・ソフトウェア}

\chaptermeta{%
  \item 10進数と2進数の相互変換ができる
  \item AND/OR/NOT/XORの論理演算を使える
  \item CPU・メモリの階層構造を説明できる
  \item RAID 0/1/5の違いを即答できる
  \item OSの役割とOSSの特徴を整理できる
}{90分}{3}{第1--4章}

テクノロジ系の最初は「コンピュータの基礎」です。
2進数・論理演算は\textbf{計算問題}として出題され、
CPU/メモリ/RAIDは\textbf{用語の識別問題}として頻出します。

% =============================================================
\section{数の表現}
% =============================================================

\begin{teigi}[基数変換の基本]
\begin{itemize}
  \item \termtag{10進数→2進数}：2で割り続け、余りを下から読む
  \item \termtag{2進数→10進数}：各桁に位の重み（$2^0, 2^1, 2^2, \ldots$）を掛けて足す
  \item \termtag{2進数→16進数}：4桁ずつ区切って16進に変換
  \item 16進数は先頭に「0x」や末尾に「H」をつけて表記（例：0x1F、1FH）
\end{itemize}
\end{teigi}

% --- 例題 ---
\begin{reidai}{10進数から2進数への変換}
10進数の13を2進数で表したものはどれか。

\sentakushi{1011}{1100}{1101}{1110}
\end{reidai}

\begin{shishin}
2で割っていく。$13 \div 2 = 6 \cdots 1$、$6 \div 2 = 3 \cdots 0$、$3 \div 2 = 1 \cdots 1$、$1 \div 2 = 0 \cdots 1$。
余りを下から読むと$1101$。
\end{shishin}

\begin{kaitou}
\textbf{正解：ウ（1101）}

$13 = 8 + 4 + 1 = 2^3 + 2^2 + 2^0 = 1101_{(2)}$。
別解として位の重みで分解する方法でも確認できる。
\end{kaitou}

\begin{minuki}[見抜き方]
\begin{itemize}
  \item 小さい数（0〜15）は暗記してしまうと速い：$1010=10, 1100=12, 1101=13, 1111=15$
  \item 2の累乗を覚える：$1, 2, 4, 8, 16, 32, 64, 128, 256, 512, 1024$
  \item 検算は2進数→10進数変換で行う
\end{itemize}
\end{minuki}

\subsection{補数と負の数}

\begin{teigi}[2の補数]
コンピュータ内部で負の数を表す方式。
\begin{enumerate}
  \item 全ビットを反転（0→1、1→0）
  \item 結果に1を足す
\end{enumerate}
例：$0101_{(2)}$（＝5）の2の補数 → $1010 + 1 = 1011_{(2)}$（＝$-5$）。
\end{teigi}

% =============================================================
\section{論理演算}
% =============================================================

\begin{hikaku}[4つの基本論理演算]
\comptable{演算 & 記号 & 結果が1になる条件}{
AND（論理積） & $\cdot$ & 両方1のときだけ \\
OR（論理和） & $+$ & どちらか一方でも1 \\
NOT（否定） & $\overline{A}$ & 入力の反転 \\
XOR（排他的論理和） & $\oplus$ & 2つが異なるとき
}
\end{hikaku}

\begin{pitfall}[ORとXORの混同]
\begin{itemize}
  \item OR：「1 OR 1 = 1」（どちらかが1ならOK）
  \item XOR：「1 XOR 1 = 0」（同じなら0、違えば1）
  \item 試験では「排他的」というキーワードでXORを識別
\end{itemize}
\end{pitfall}

% =============================================================
\section{アルゴリズムの基礎}
% =============================================================

\begin{teigi}[基本的なアルゴリズム]
\begin{itemize}
  \item \termtag{線形探索}：先頭から順に探す（単純だがデータ量に比例して遅い）
  \item \termtag{2分探索}：ソート済みデータを半分ずつ絞る（高速：$O(\log n)$）
  \item \termtag{バブルソート}：隣り合うデータを比較・交換して並べ替え
  \item \termtag{選択ソート}：最小値を見つけて先頭と交換を繰り返す
\end{itemize}
\end{teigi}

\begin{pitfall}[「2分探索はソート不要」は誤り]
2分探索の前提は\textbf{データがソート済み}であること。
未整列データに2分探索を適用しても正しい結果は得られない。
\end{pitfall}

% =============================================================
\section{CPU・メモリの階層構造}
% =============================================================

\begin{teigi}[記憶装置の階層]
速度が速い順（＝容量が小さい順）：
\begin{enumerate}
  \item \termtag{レジスタ}：CPU内部の超高速記憶（数十〜数百個）
  \item \termtag{キャッシュメモリ}：CPUとメインメモリの間の高速バッファ（SRAM）
  \item \termtag{メインメモリ（主記憶）}：実行中プログラムが置かれる（DRAM）
  \item \termtag{SSD/HDD（補助記憶）}：電源を切ってもデータが残る
\end{enumerate}
上に行くほど\textbf{高速・小容量・高価格}、下に行くほど\textbf{低速・大容量・低価格}。
\end{teigi}

\begin{hikaku}[SRAMとDRAM]
\comptable{特性 & SRAM & DRAM}{
速度 & 高速 & 低速 \\
容量あたりの価格 & 高い & 安い \\
リフレッシュ & 不要 & 必要 \\
用途 & キャッシュメモリ & メインメモリ
}
\end{hikaku}

\subsection{CPUの性能指標}

\begin{teigi}[クロック周波数とCPI]
\begin{itemize}
  \item \termtag{クロック周波数}：CPUが1秒間に刻むクロック数（単位：Hz）。大きいほど高速。
  \item \termtag{MIPS}（Million Instructions Per Second）：1秒間に実行できる命令数（百万単位）
  \item クロック周波数1GHz、1命令に2クロック必要 → MIPS $= 1{,}000 \div 2 = 500$ MIPS
\end{itemize}
\end{teigi}

% =============================================================
\section{RAID}
% =============================================================

\begin{hikaku}[RAID 0/1/5の比較]
\comptable{RAID & 方式 & 特徴}{
RAID 0 & ストライピング & 高速だが冗長性なし（1台故障で全滅） \\
RAID 1 & ミラーリング & 同じデータを2台に書く（容量半減） \\
RAID 5 & パリティ分散 & 3台以上、1台故障まで復旧可能
}
\end{hikaku}

\begin{pitfall}[RAID 5の「1台まで」を忘れる]
RAID 5は\textbf{1台}の故障まで復旧可能。2台同時故障ではデータが失われる。
2台故障に耐えるのはRAID 6（パリティを二重化）。
\end{pitfall}

% =============================================================
\section{OS・OSS}
% =============================================================

\begin{teigi}[OSの役割]
\begin{itemize}
  \item \termtag{OS}（Operating System）：ハードウェアを抽象化し、アプリケーションに共通の実行環境を提供
  \item 主な機能：プロセス管理、メモリ管理、ファイル管理、入出力管理
  \item \termtag{仮想記憶}：主記憶が不足したとき、補助記憶の一部を主記憶のように使う仕組み
\end{itemize}
\end{teigi}

\begin{teigi}[OSSの特徴]
\termtag{OSS}（Open Source Software）：ソースコードが公開され、自由に利用・改変・再配布できるソフトウェア。
\begin{itemize}
  \item 無償とは限らない（サポートは有償の場合あり）
  \item 代表例：Linux、Apache、MySQL、PostgreSQL、Python、Firefox
  \item ライセンス例：GPL（改変後も公開義務）、MIT（制約が少ない）
\end{itemize}
\end{teigi}

\begin{pitfall}[「OSSは無料」は不正確]
OSSは\textbf{ソースコードが公開}されていることが本質。
商用サポートを有償で提供することは可能であり、OSS＝無料ではない。
\end{pitfall}

% =============================================================
\section{過去問ドリル}
% =============================================================

\shoumatsuhead{過去問ドリル：基礎理論とハードウェア}

\begin{itpmon}[\sourcebadge{original}\enspace\diffbadge{標}]{%
RAID 5の説明として、最も適切なものはどれか。}
\sentakushi
  {データを複数のディスクに分散して書き込み、高速化するが冗長性はない。}
  {同一データを2台のディスクに書き込み、1台が故障してもデータを保持できる。}
  {3台以上のディスクにデータとパリティを分散し、1台の故障まで復旧可能である。}
  {2台のディスクにそれぞれ独立したパリティを持ち、2台の故障まで復旧可能である。}
\end{itpmon}
\begin{kaisetsu}{ウ}
\textbf{RAID 5}はデータとパリティを3台以上に分散記録する。アはRAID 0、イはRAID 1、エはRAID 6の説明。
\end{kaisetsu}

\begin{itpmon}[\sourcebadge{original}\enspace\diffbadge{標}]{%
CPUのクロック周波数に関する説明として、最も適切なものはどれか。}
\sentakushi
  {クロック周波数が高いほど、1秒あたりに実行できる命令数が多くなる。}
  {クロック周波数が高いほど、記憶装置の容量が大きくなる。}
  {クロック周波数は、ネットワークの通信速度を表す指標である。}
  {クロック周波数は、ディスプレイの解像度を決定する要素である。}
\end{itpmon}
\begin{kaisetsu}{ア}
\textbf{クロック周波数}はCPUの動作速度の指標。周波数が高いほど単位時間に多くの命令を処理できる。記憶容量・通信速度・解像度とは直接関係しない。
\end{kaisetsu}

\begin{itpmon}[\sourcebadge{original}\enspace\diffbadge{標}]{%
OSSに関する説明として、最も適切なものはどれか。}
\sentakushi
  {ソースコードが非公開のソフトウェアである。}
  {利用にはメーカーのライセンス契約が必須である。}
  {ソースコードが公開され、自由に利用・改変・再配布が認められている。}
  {企業が開発したソフトウェアの体験版である。}
\end{itpmon}
\begin{kaisetsu}{ウ}
\textbf{OSS}はソースコードの公開が本質。GPL、MIT、Apacheライセンスなど、様々なライセンス形態がある。自由に利用・改変・再配布が認められている。
\end{kaisetsu}

% =============================================================
\section{タイムアタック}
% =============================================================

\begin{timeattack}{基礎理論クイックチェック}{6分}{5問中4問}

\begin{itpmon}[\diffbadge{基}]{10進数の10を2進数で表したものはどれか。}
\sentakushi{1000}{1001}{1010}{1011}
\end{itpmon}
\begin{kaisetsu}{ウ}
$10 = 8 + 2 = 2^3 + 2^1 = 1010_{(2)}$。
\end{kaisetsu}

\begin{itpmon}[\diffbadge{基}]{論理演算「1 AND 0」の結果はどれか。}
\sentakushi{0}{1}{2}{不定}
\end{itpmon}
\begin{kaisetsu}{ア}
AND（論理積）は両方1のときだけ1。$1 \text{ AND } 0 = 0$。
\end{kaisetsu}

\begin{itpmon}[\diffbadge{標}]{キャッシュメモリの説明として正しいものはどれか。}
\sentakushi
  {電源を切ってもデータが残る記憶装置}
  {CPUとメインメモリの間に置かれる高速な記憶装置}
  {大容量で低価格な磁気ディスク装置}
  {プログラムの実行中にデータを一時的に退避する補助記憶装置}
\end{itpmon}
\begin{kaisetsu}{イ}
\textbf{キャッシュメモリ}はCPUとメインメモリの速度差を埋めるSRAM。アはSSD/HDD、ウもHDD、エはスワップ領域の説明。
\end{kaisetsu}

\begin{itpmon}[\diffbadge{基}]{RAID 1の方式名はどれか。}
\sentakushi{ストライピング}{ミラーリング}{パリティ分散}{二重パリティ}
\end{itpmon}
\begin{kaisetsu}{イ}
RAID 1 = \textbf{ミラーリング}。アはRAID 0、ウはRAID 5、エはRAID 6。
\end{kaisetsu}

\begin{itpmon}[\diffbadge{標}]{仮想記憶の説明として正しいものはどれか。}
\sentakushi
  {CPUの内部にある超高速な記憶領域}
  {主記憶の不足分を補助記憶で補い、見かけ上の主記憶を拡大する仕組み}
  {データを複数のディスクに分散する技術}
  {ネットワーク上の記憶領域を共有する仕組み}
\end{itpmon}
\begin{kaisetsu}{イ}
\textbf{仮想記憶}は主記憶が不足したとき、補助記憶（SSD/HDD）の一部を使って見かけ上の主記憶容量を拡大する技術。
\end{kaisetsu}

\end{timeattack}

% =============================================================
\section{よくあるミスと到達確認}
% =============================================================

\begin{yokumiss}[この章でよくあるミス]
\begin{enumerate}[leftmargin=1.6em, itemsep=5pt]
  \item \textbf{2進数変換で余りの読む方向を間違える}——下から上に読む。
  \item \textbf{ORとXORを混同する}——OR: 1 OR 1 = 1、XOR: 1 XOR 1 = 0。
  \item \textbf{SRAMとDRAMの用途を逆にする}——SRAM=キャッシュ、DRAM=メインメモリ。
  \item \textbf{RAID 0に冗長性があると思い込む}——RAID 0は高速化のみ、冗長性なし。
  \item \textbf{OSSを「無料」と同義に扱う}——本質はソースコード公開。商用サポートは有償可。
\end{enumerate}
\end{yokumiss}

\begin{tatsaku}
  \item 10進数と2進数を相互変換できる
  \item AND/OR/NOT/XORの真理値表を書ける
  \item 記憶装置の階層（レジスタ→キャッシュ→メインメモリ→補助記憶）を言える
  \item SRAMとDRAMの違いを3点以上挙げられる
  \item RAID 0/1/5の違いを即答できる
  \item 仮想記憶の仕組みを説明できる
  \item OSSの定義と代表例を3つ以上言える
  \item GPLとMITライセンスの違いを説明できる
\end{tatsaku}

\nextchapter{%
ネットワークとデータベースに進みます。TCP/IP、プロトコル、SQL、正規化——
「つながる仕組み」と「データの管理」を学びます。
}

% =============================================================
% 06_network.tex — ネットワークとデータベース
% =============================================================
\chapter{ネットワークとデータベース}

\chaptermeta{%
  \item TCP/IPの4層モデルを説明できる
  \item 主要プロトコルの役割を即答できる
  \item SQLの基本構文（SELECT/WHERE/ORDER BY）が読める
  \item 正規化の目的と第1〜第3正規形を区別できる
}{90分}{3}{第5章}

テクノロジ系の中盤は「つながる仕組み」と「データの管理」です。
TCP/IPの階層は\textbf{プロトコルの分類問題}で頻出し、
SQLは\textbf{文法を読めるかどうか}がそのまま得点に直結します。

% =============================================================
\section{TCP/IPの4層モデル}
% =============================================================

\begin{teigi}[TCP/IP 4層]
\begin{enumerate}
  \item \termtag{アプリケーション層}：HTTP, HTTPS, SMTP, POP3, FTP, DNS, DHCP, NTP
  \item \termtag{トランスポート層}：TCP, UDP
  \item \termtag{インターネット層}：IP, ICMP, ARP
  \item \termtag{ネットワークインターフェース層}：Ethernet, Wi-Fi
\end{enumerate}
OSI参照モデル（7層）との対応も問われるが、Iパスでは4層が中心。
\end{teigi}

\begin{hikaku}[主要プロトコルの整理]
\comptable{プロトコル & 層 & 役割}{
HTTP/HTTPS & アプリケーション & Webページの送受信（Sは暗号化あり） \\
SMTP & アプリケーション & メール送信 \\
POP3/IMAP & アプリケーション & メール受信（POP3=DL型、IMAP=サーバ参照型） \\
FTP & アプリケーション & ファイル転送 \\
DNS & アプリケーション & ドメイン名 ↔ IPアドレスの変換 \\
DHCP & アプリケーション & IPアドレスの自動割り当て \\
NTP & アプリケーション & 時刻同期 \\
TCP & トランスポート & 信頼性のある通信（コネクション型） \\
UDP & トランスポート & 高速だが信頼性なし（コネクションレス型）
}
\end{hikaku}

% =============================================================
\section{TCPとUDP}
% =============================================================

\begin{hikaku}[TCP vs UDP]
\comptable{特性 & TCP & UDP}{
通信方式 & コネクション型 & コネクションレス型 \\
信頼性 & 高い（再送制御あり） & 低い（再送なし） \\
速度 & 相対的に遅い & 高速 \\
用途 & Web・メール・ファイル転送 & 動画配信・VoIP・DNS問い合わせ
}
\end{hikaku}

% --- 例題 ---
\begin{reidai}{TCPとUDPの識別}
リアルタイムの動画配信に適したトランスポート層のプロトコルはどれか。
信頼性よりも速度を重視する場合を考えよ。

\sentakushi{FTP}{HTTP}{TCP}{UDP}
\end{reidai}

\begin{shishin}
「リアルタイム」「速度重視」がキーワード。FTPとHTTPはアプリケーション層で、トランスポート層ではない。
TCPは信頼性重視で再送制御がある分、遅延が発生しやすい。
UDPは再送せず高速だが、データが欠落することがある。
動画配信では多少の欠落よりリアルタイム性を優先する。
\end{shishin}

\begin{kaitou}
\textbf{正解：エ（UDP）}

UDPはコネクションレス型で再送制御がなく高速。
リアルタイム配信・VoIP・オンラインゲームなど、遅延より速度を重視する場面で使われる。
\end{kaitou}

% =============================================================
\section{DNS・DHCP}
% =============================================================

\begin{teigi}[DNSとDHCP]
\begin{itemize}
  \item \termtag{DNS}（Domain Name System）：ドメイン名（www.example.com）をIPアドレス（192.168.1.1）に変換する「電話帳」
  \item \termtag{DHCP}（Dynamic Host Configuration Protocol）：ネットワーク接続時にIPアドレスを自動的に割り当てる
\end{itemize}
\end{teigi}

\begin{pitfall}[DNSとDHCPの混同]
\begin{itemize}
  \item DNS → \textbf{名前解決}（ドメイン名 ↔ IPアドレスの変換）
  \item DHCP → \textbf{アドレス配布}（IPアドレスの自動割り当て）
  \item 「DNSは翻訳、DHCPは配布」と覚える
\end{itemize}
\end{pitfall}

% =============================================================
\section{サブネットマスク}
% =============================================================

\begin{teigi}[IPアドレスとサブネットマスク]
\begin{itemize}
  \item \termtag{IPアドレス}：ネットワーク上の機器を識別する番号（IPv4は32ビット）
  \item \termtag{サブネットマスク}：IPアドレスの「ネットワーク部」と「ホスト部」の境界を示す
  \item 例：255.255.255.0 → 上位24ビットがネットワーク部、下位8ビットがホスト部
  \item 同じネットワーク部を持つ機器同士が「同一ネットワーク」
\end{itemize}
\end{teigi}

\begin{pitfall}[サブネットマスクの計算ミス]
サブネットマスク255.255.255.0の場合、ホスト部は8ビット。
使えるアドレス数は $2^8 - 2 = 254$（ネットワークアドレスとブロードキャストを除く）。
「$-2$」を忘れがち。
\end{pitfall}

% =============================================================
\section{RDB（関係データベース）の基礎}
% =============================================================

\begin{teigi}[RDBの基本用語]
\begin{itemize}
  \item \termtag{テーブル（表）}：データを行と列で格納する
  \item \termtag{レコード（行/タプル）}：1件分のデータ
  \item \termtag{フィールド（列/属性）}：データの項目
  \item \termtag{主キー}：各レコードを一意に識別するフィールド（NULLや重複不可）
  \item \termtag{外部キー}：他のテーブルの主キーを参照するフィールド（参照整合性の確保）
\end{itemize}
\end{teigi}

\begin{hikaku}[主キーと外部キー]
\comptable{項目 & 主キー & 外部キー}{
役割 & 自テーブルのレコードを一意に識別 & 他テーブルの主キーを参照 \\
NULL & 不可 & 許可される場合あり \\
重複 & 不可 & 許可される場合あり \\
目的 & レコードの特定 & テーブル間の関連づけ
}
\end{hikaku}

% =============================================================
\section{SQL}
% =============================================================

\begin{teigi}[SQLの基本構文]
\begin{itemize}
  \item \termtag{SELECT}：取得する列を指定
  \item \termtag{FROM}：対象テーブルを指定
  \item \termtag{WHERE}：条件で絞り込み
  \item \termtag{ORDER BY}：並び替え（ASC=昇順、DESC=降順）
  \item \termtag{GROUP BY}：グループ化（集計関数 COUNT, SUM, AVG, MAX, MIN と併用）
  \item \termtag{HAVING}：GROUP BYの結果に対する条件
\end{itemize}
\end{teigi}

\begin{advice}[SQLは「読めればOK」]
Iパスでは SQL を書く問題は出ず、\textbf{SELECT文を読んで結果を選ぶ}形式。
WHERE の条件と ORDER BY の並び順を追えれば正解できる。
\end{advice}

% =============================================================
\section{正規化}
% =============================================================

\begin{teigi}[正規化の3段階]
\begin{itemize}
  \item \termtag{第1正規形}：繰り返し項目を排除（1つのセルに1つの値）
  \item \termtag{第2正規形}：部分関数従属を排除（主キーの一部だけに依存する列を別テーブルに分離）
  \item \termtag{第3正規形}：推移的関数従属を排除（主キー以外の列に依存する列を別テーブルに分離）
\end{itemize}
正規化の目的は\textbf{データの冗長性を排除}し、\textbf{更新時の矛盾を防ぐ}こと。
\end{teigi}

\begin{pitfall}[正規化のレベルを混同]
\begin{itemize}
  \item 第1：「繰り返しがない」だけではまだ第1
  \item 第2：「主キー全体に依存」＝部分関数従属がない
  \item 第3：「主キー以外への依存もない」＝推移的関数従属がない
  \item 「1→繰り返し除去、2→部分除去、3→推移除去」の語呂で覚える
\end{itemize}
\end{pitfall}

% =============================================================
\section{過去問ドリル}
% =============================================================

\shoumatsuhead{過去問ドリル：ネットワークとDB}

\begin{itpmon}[\sourcebadge{original}\enspace\diffbadge{標}]{%
次のSQL文を実行した結果として得られるものはどれか。

\texttt{SELECT 氏名, 点数 FROM 成績表 WHERE 点数 >= 80 ORDER BY 点数 DESC}}
\sentakushi
  {成績表から点数が80以上のレコードを点数の昇順で取得する。}
  {成績表から点数が80以上のレコードを点数の降順で取得する。}
  {成績表から点数が80未満のレコードを点数の昇順で取得する。}
  {成績表からすべてのレコードを点数の降順で取得する。}
\end{itpmon}
\begin{kaisetsu}{イ}
WHERE 点数 >= 80 で80点以上を絞り込み、ORDER BY 点数 \textbf{DESC}（降順）で並べ替え。DESCは大きい順、ASCは小さい順。
\end{kaisetsu}

\begin{itpmon}[\sourcebadge{original}\enspace\diffbadge{標}]{%
RDB（関係データベース）における主キーと外部キーの説明として、最も適切なものはどれか。}
\sentakushi
  {主キーはNULLを許容し、外部キーはNULLを許容しない。}
  {主キーはレコードを一意に識別し、外部キーは他テーブルの主キーを参照する。}
  {主キーも外部キーも、同一テーブル内でのみ使用される。}
  {外部キーは主キーの代わりに使うことができる。}
\end{itpmon}
\begin{kaisetsu}{イ}
\textbf{主キー}はレコードの一意識別（NULL・重複不可）、\textbf{外部キー}は他テーブルの主キーを参照してテーブル間を関連づける。アは逆。
\end{kaisetsu}

\begin{itpmon}[\sourcebadge{original}\enspace\diffbadge{強}]{%
サブネットマスクが255.255.255.0のネットワークで、ホストに割り当てられるIPアドレスの最大数はどれか。}
\sentakushi
  {252}
  {254}
  {255}
  {256}
\end{itpmon}
\begin{kaisetsu}{イ}
ホスト部は8ビットで $2^8 = 256$ 通りだが、ネットワークアドレス（全0）とブロードキャストアドレス（全1）を除くため $256 - 2 = \textbf{254}$。
\end{kaisetsu}

% =============================================================
\section{よくあるミスと到達確認}
% =============================================================

\begin{yokumiss}[この章でよくあるミス]
\begin{enumerate}[leftmargin=1.6em, itemsep=5pt]
  \item \textbf{TCPとUDPを逆にする}——TCP=信頼性重視、UDP=速度重視。「Tは丁寧（Teinei）」で覚える。
  \item \textbf{DNSとDHCPを混同する}——DNS=名前解決、DHCP=アドレス自動配布。
  \item \textbf{SQLのDESC/ASCを逆に覚える}——DESC=降順（大→小）、ASC=昇順（小→大）。
  \item \textbf{主キーにNULLが許される}と思い込む——主キーはNULL禁止・重複禁止。
  \item \textbf{サブネットマスクの計算で$-2$を忘れる}——ネットワークアドレスとブロードキャストは除外。
  \item \textbf{正規化のレベルを混同する}——1=繰り返し除去、2=部分従属除去、3=推移従属除去。
\end{enumerate}
\end{yokumiss}

\begin{tatsaku}
  \item TCP/IP 4層と各層のプロトコルを3つ以上言える
  \item TCPとUDPの違いを3点挙げられる
  \item DNS/DHCPの役割を即答できる
  \item サブネットマスクからホスト数を計算できる
  \item SQLのSELECT文を読んで結果を予測できる
  \item 主キーと外部キーの違いを説明できる
  \item 正規化の第1〜第3正規形を区別できる
  \item POP3とIMAPの違いを説明できる
\end{tatsaku}

\nextchapter{%
セキュリティとAIに進みます。CIA三要素、暗号方式、マルウェア対策、そしてAI/ML/DL——
テクノロジ系の最重要テーマを攻略します。
}

% =============================================================
% 07_security.tex — セキュリティとAI
% =============================================================
\chapter{セキュリティとAI}

\chaptermeta{%
  \item CIA三要素を具体例と共に説明できる
  \item 認証の3要素を区別できる
  \item 公開鍵暗号とデジタル署名の仕組みを理解する
  \item 主要な脅威と対策を整理できる
  \item AI/ML/DLの関係と基本用語を説明できる
}{90分}{3}{第5--6章}

テクノロジ系で最も出題数が多いのがセキュリティ分野です。
近年はAI関連の出題も増加傾向にあります。
\textbf{暗号方式の仕組み}と\textbf{攻撃手法の名前}を確実に押さえましょう。

% =============================================================
\section{CIA三要素}
% =============================================================

\begin{teigi}[情報セキュリティのCIA]
\begin{itemize}
  \item \termtag{機密性}（Confidentiality）：許可された人だけがアクセスできること
  \item \termtag{完全性}（Integrity）：情報が改ざん・破壊されていないこと
  \item \termtag{可用性}（Availability）：必要なときに利用できること
\end{itemize}
この3つのバランスを保つことが情報セキュリティの目的。
\end{teigi}

\begin{hikaku}[CIA三要素の具体例]
\comptable{要素 & 脅威の例 & 対策の例}{
機密性 & 不正アクセス・盗聴 & 暗号化・アクセス制御 \\
完全性 & データ改ざん・ウイルス & ハッシュ値検証・デジタル署名 \\
可用性 & DoS攻撃・機器故障 & 冗長化・バックアップ
}
\end{hikaku}

% =============================================================
\section{認証の3要素}
% =============================================================

\begin{teigi}[認証の3要素]
\begin{itemize}
  \item \termtag{知識情報}：本人だけが知っていること（パスワード、PINコード、秘密の質問）
  \item \termtag{所持情報}：本人だけが持っているもの（ICカード、スマートフォン、ワンタイムパスワード生成器）
  \item \termtag{生体情報}：本人の身体的特徴（指紋、顔、虹彩、静脈パターン）
\end{itemize}
\termtag{二要素認証}は上記3種から\textbf{異なる2種}を組み合わせること。
パスワード＋秘密の質問は両方「知識」なので二要素認証にはならない。
\end{teigi}

\begin{pitfall}[「二段階認証」＝「二要素認証」ではない]
\begin{itemize}
  \item 二段階認証：認証を2回行うこと（同じ要素でもよい）
  \item 二要素認証：異なる種類の要素を2つ組み合わせること
  \item パスワード＋SMS認証 → 知識＋所持 → 二要素認証
  \item パスワード＋秘密の質問 → 知識＋知識 → 二段階だが一要素
\end{itemize}
\end{pitfall}

% =============================================================
\section{公開鍵暗号とデジタル署名}
% =============================================================

\begin{hikaku}[共通鍵暗号と公開鍵暗号]
\comptable{項目 & 共通鍵暗号 & 公開鍵暗号}{
鍵の数 & 送受信者で同じ1つの鍵 & 公開鍵と秘密鍵のペア \\
速度 & 高速 & 低速 \\
鍵配送問題 & あり（安全に鍵を渡す必要） & なし（公開鍵は公開してOK） \\
代表例 & AES & RSA
}
\end{hikaku}

\begin{teigi}[デジタル署名の仕組み]
\begin{enumerate}
  \item 送信者がメッセージの\textbf{ハッシュ値}を計算
  \item 送信者の\textbf{秘密鍵}でハッシュ値を暗号化 → 署名
  \item 受信者は送信者の\textbf{公開鍵}で署名を復号
  \item 復号したハッシュ値と、自分で計算したハッシュ値を比較
  \item 一致すれば\textbf{改ざんなし}かつ\textbf{本人が送信}と確認
\end{enumerate}
公開鍵暗号の「逆使い」——暗号化は受信者の公開鍵、署名は送信者の秘密鍵。
\end{teigi}

% --- 例題 ---
\begin{reidai}{公開鍵暗号 vs デジタル署名}
公開鍵暗号方式による暗号化通信とデジタル署名について、正しい組合せはどれか。

\begin{center}
\small
\begin{tabular}{lll}
\toprule
 & 暗号化通信 & デジタル署名 \\
\midrule
ア & 受信者の公開鍵で暗号化 & 送信者の公開鍵で署名 \\
イ & 受信者の公開鍵で暗号化 & 送信者の秘密鍵で署名 \\
ウ & 送信者の秘密鍵で暗号化 & 受信者の公開鍵で署名 \\
エ & 送信者の公開鍵で暗号化 & 受信者の秘密鍵で署名 \\
\bottomrule
\end{tabular}
\end{center}

\sentakushi{ア}{イ}{ウ}{エ}
\end{reidai}

\begin{shishin}
暗号化通信：「相手しか復号できない」→ 受信者の公開鍵で暗号化、受信者の秘密鍵で復号。
デジタル署名：「本人が送った証拠」→ 送信者の秘密鍵で署名、送信者の公開鍵で検証。
\end{shishin}

\begin{kaitou}
\textbf{正解：イ}

暗号化通信 = 受信者の公開鍵で暗号化。デジタル署名 = 送信者の秘密鍵で署名。
「暗号は相手の公開鍵、署名は自分の秘密鍵」と覚える。
\end{kaitou}

\begin{minuki}[見抜き方]
\begin{itemize}
  \item 暗号化：「受信者の公開鍵」→ 受信者だけが復号できる
  \item デジタル署名：「送信者の秘密鍵」→ 送信者本人の証明
  \item 「暗号は受信者、署名は送信者」——主役が入れ替わる
\end{itemize}
\end{minuki}

% =============================================================
\section{脅威と対策}
% =============================================================

\begin{hikaku}[主要なマルウェアと攻撃手法]
\comptable{種類 & 特徴 & 対策}{
ランサムウェア & データを暗号化して身代金を要求 & バックアップ・OS更新 \\
フィッシング & 偽サイトで情報を騙し取る & URL確認・二要素認証 \\
標的型攻撃 & 特定組織を狙った巧妙なメール & 訓練・メールフィルタ \\
DoS/DDoS & 大量通信でサービスを停止させる & CDN・IP制限 \\
SQLインジェクション & 入力フォームからDB操作 & 入力値の無害化 \\
クロスサイトスクリプティング & Webページに悪意のスクリプト注入 & エスケープ処理 \\
ソーシャルエンジニアリング & 人間の心理を利用して情報を入手 & セキュリティ教育
}
\end{hikaku}

\begin{pitfall}[ランサムウェアの対処で「身代金を払う」を選びがち]
\begin{itemize}
  \item 身代金を払ってもデータが復旧する保証はない
  \item 正しい対策は\textbf{日頃からバックアップ}を取ること
  \item 感染時は即座にネットワークから切り離す
\end{itemize}
\end{pitfall}

% =============================================================
\section{信頼性設計}
% =============================================================

\begin{hikaku}[信頼性設計の比較]
\comptable{設計思想 & 方針 & 具体例}{
フォールトトレランス & 故障しても動き続ける & 冗長化・RAID \\
フェールセーフ & 故障時に安全側に倒す & 踏切の遮断機（故障→閉じる） \\
フェールソフト & 故障時に機能を縮退して継続 & 空調故障→換気だけ継続 \\
フールプルーフ & 人間のミスを防ぐ設計 & コンセントの形状・確認ダイアログ
}
\end{hikaku}

% =============================================================
\section{AI概論}
% =============================================================

\begin{teigi}[AI / ML / DLの関係]
\begin{itemize}
  \item \termtag{AI}（人工知能）：人間の知能を模倣する技術の総称
  \item \termtag{ML}（機械学習）：AIの一分野。データからパターンを自動学習
  \item \termtag{DL}（深層学習/ディープラーニング）：MLの一手法。多層ニューラルネットワーク
  \item 関係：AI $\supset$ ML $\supset$ DL（入れ子構造）
\end{itemize}
\end{teigi}

\begin{hikaku}[機械学習の3分類]
\comptable{学習方式 & 特徴 & 例}{
教師あり学習 & 正解ラベル付きデータで学習 & 画像分類・スパム判定 \\
教師なし学習 & 正解なしでパターンを発見 & クラスタリング・次元削減 \\
強化学習 & 報酬を最大化する行動を学習 & ゲームAI・ロボット制御
}
\end{hikaku}

\begin{teigi}[AI関連の重要用語]
\begin{itemize}
  \item \termtag{AIバイアス}：学習データの偏りにより、AIが不公平な判断をすること
  \item \termtag{説明可能AI（XAI）}：AIの判断根拠を人間が理解できるようにする技術
  \item \termtag{生成AI}：テキスト・画像などのコンテンツを生成するAI（例：ChatGPT）
  \item \termtag{過学習}：訓練データに適合しすぎて、未知データへの予測精度が下がる現象
\end{itemize}
\end{teigi}

% =============================================================
\section{過去問ドリル}
% =============================================================

\shoumatsuhead{過去問ドリル：セキュリティとAI}

\begin{itpmon}[\sourcebadge{R7}\enspace\diffbadge{標}]{%
JIS Q 27017に基づく、クラウドサービス利用におけるセキュリティ管理策の説明として、最も適切なものはどれか。}
\sentakushi
  {クラウドサービス利用者が独自にセキュリティ基準を策定する指針}
  {クラウドサービスに関する情報セキュリティ管理策の実践規範}
  {オンプレミス環境のみを対象とした情報セキュリティ規格}
  {個人情報保護に特化したクラウドセキュリティのガイドライン}
\end{itpmon}
\begin{kaisetsu}{イ}
\textbf{ISO/IEC 27017}（JIS Q 27017）はクラウドサービスに関する情報セキュリティ管理策の実践規範。ISO 27001/27002を補完するクラウド固有のガイドライン。
\end{kaisetsu}

\begin{itpmon}[\sourcebadge{R7}\enspace\diffbadge{標}]{%
AIが学習データの偏りによって、特定の属性を持つ人々に対して不公平な判断を行ってしまう問題を何というか。}
\sentakushi
  {過学習}
  {AIバイアス}
  {シンギュラリティ}
  {チューリングテスト}
\end{itpmon}
\begin{kaisetsu}{イ}
\textbf{AIバイアス}は学習データの偏り（性別・人種等）によりAIが不公平な判断をすること。アは訓練データへの過剰適合、ウはAIが人間の知能を超える仮説上の時点、エはAIの知能を判定するテスト。
\end{kaisetsu}

\begin{itpmon}[\sourcebadge{original}\enspace\diffbadge{標}]{%
ランサムウェアに関する説明として、最も適切なものはどれか。}
\sentakushi
  {Webサイトを閲覧しただけで感染し、広告を大量に表示するソフトウェア}
  {コンピュータ内のファイルを暗号化し、復号の見返りに金銭を要求するマルウェア}
  {キーボードの入力内容を記録して外部に送信するソフトウェア}
  {正常なソフトウェアを装ってインストールさせ、裏で不正な動作を行うマルウェア}
\end{itpmon}
\begin{kaisetsu}{イ}
\textbf{ランサムウェア}（ransom=身代金）はファイルを暗号化し、復号と引き換えに身代金を要求する。アはアドウェア、ウはキーロガー、エはトロイの木馬。
\end{kaisetsu}

% =============================================================
\section{タイムアタック}
% =============================================================

\begin{timeattack}{セキュリティ・AIクイックチェック}{6分}{5問中4問}

\begin{itpmon}[\diffbadge{基}]{情報セキュリティのCIA三要素に含まれないものはどれか。}
\sentakushi{機密性}{完全性}{可用性}{効率性}
\end{itpmon}
\begin{kaisetsu}{エ}
CIAは機密性（C）・完全性（I）・可用性（A）。\textbf{効率性}はCIAに含まれない。
\end{kaisetsu}

\begin{itpmon}[\diffbadge{標}]{デジタル署名で、送信者が署名に使うのはどの鍵か。}
\sentakushi{受信者の公開鍵}{受信者の秘密鍵}{送信者の公開鍵}{送信者の秘密鍵}
\end{itpmon}
\begin{kaisetsu}{エ}
デジタル署名は\textbf{送信者の秘密鍵}で署名し、送信者の公開鍵で検証する。
\end{kaisetsu}

\begin{itpmon}[\diffbadge{基}]{フェールセーフの例として最も適切なものはどれか。}
\sentakushi
  {システムの一部が故障しても全体は動き続ける}
  {故障時に安全な状態で停止する}
  {人間がミスをしても危険な動作にならない設計}
  {機能を縮退して運転を継続する}
\end{itpmon}
\begin{kaisetsu}{イ}
\textbf{フェールセーフ}は故障時に安全側に倒す設計。アはフォールトトレランス、ウはフールプルーフ、エはフェールソフト。
\end{kaisetsu}

\begin{itpmon}[\diffbadge{標}]{機械学習で、正解ラベル付きデータを使って学習する方式はどれか。}
\sentakushi{教師あり学習}{教師なし学習}{強化学習}{転移学習}
\end{itpmon}
\begin{kaisetsu}{ア}
\textbf{教師あり学習}は正解ラベル付きデータで学習。イは正解なし、ウは報酬最大化、エは別タスクで学習した知識を転用する手法。
\end{kaisetsu}

\begin{itpmon}[\diffbadge{標}]{SQLインジェクション攻撃の対策として最も適切なものはどれか。}
\sentakushi
  {強力なパスワードを設定する}
  {入力値の無害化（サニタイジング）を行う}
  {ファイアウォールを設置する}
  {通信を暗号化する}
\end{itpmon}
\begin{kaisetsu}{イ}
SQLインジェクションはWebフォームの入力値にSQL文を混入させる攻撃。対策は\textbf{入力値の無害化}（プレースホルダやエスケープ処理）。
\end{kaisetsu}

\end{timeattack}

% =============================================================
\section{よくあるミスと到達確認}
% =============================================================

\begin{yokumiss}[この章でよくあるミス]
\begin{enumerate}[leftmargin=1.6em, itemsep=5pt]
  \item \textbf{公開鍵暗号で「送信者の公開鍵で暗号化」と答える}——暗号化は受信者の公開鍵。
  \item \textbf{デジタル署名で「公開鍵で署名」と答える}——署名は送信者の秘密鍵。
  \item \textbf{フェールセーフとフールプルーフを混同する}——セーフは「故障→安全停止」、プルーフは「人為ミス防止」。
  \item \textbf{二段階認証を二要素認証と同一視する}——異なる種類の要素を使うかどうかが違い。
  \item \textbf{AIバイアスと過学習を混同する}——バイアスはデータの偏り、過学習は訓練データへの過剰適合。
\end{enumerate}
\end{yokumiss}

\begin{tatsaku}
  \item CIA三要素を具体例付きで説明できる
  \item 認証の3要素をそれぞれ2つ以上例示できる
  \item 公開鍵暗号とデジタル署名の鍵の使い分けを即答できる
  \item ランサムウェア・フィッシング・標的型攻撃を区別できる
  \item フォールトトレランス/フェールセーフ/フェールソフト/フールプルーフを区別できる
  \item AI $\supset$ ML $\supset$ DLの関係を説明できる
  \item 教師あり/教師なし/強化学習を区別できる
  \item AIバイアスと説明可能AIの意味を説明できる
\end{tatsaku}

\nextchapter{%
分野横断の整理と試験戦略に進みます。全章の知識をつなぎ、
ひっかけの見抜き方、本試験の時間配分を攻略します。
}


% ── PART IV: 総仕上げ ──
\partpage{IV}{総仕上げ}
% =============================================================
% 08_crossdomain.tex — 分野横断の整理と試験戦略
% =============================================================
\chapter{分野横断の整理と試験戦略}

\chaptermeta{%
  \item 横断テーマ（法律/セキュリティ/マネジメント/AI）を整理できる
  \item ひっかけ選択肢のパターンを見抜ける
  \item 本試験の時間配分を設計できる
  \item 捨て問の判断基準を持てる
}{45分}{3}{第1--7章すべて}

本章は知識のインプットではなく「戦略」の章です。
7章までに学んだ内容を\textbf{横断的に整理}し、
本試験で確実に合格点を取る\textbf{戦術}を身につけます。

% =============================================================
\section{横断テーマ4つ}
% =============================================================

\begin{teigi}[分野をまたぐ4大テーマ]
Iパスの出題は3分野に分かれますが、次の4テーマは分野を横断して出題されます。
\begin{enumerate}
  \item \termtag{法律・制度}：個人情報保護法、著作権法、不正アクセス禁止法、労働者派遣法
  \item \termtag{セキュリティ}：CIA三要素、暗号、認証（ストラテジ系でも出る）
  \item \termtag{マネジメント手法}：PDCAサイクル、リスク管理（全分野共通）
  \item \termtag{AI/データ活用}：AI、ビッグデータ、IoT（全分野から出題）
\end{enumerate}
\end{teigi}

\begin{hikaku}[横断テーマ：どの章で学んだか]
\comptable{テーマ & 主に学んだ章 & 他の章での出題例}{
法律・制度 & 第2章（法務） & 第7章（不正アクセス禁止法） \\
セキュリティ & 第7章 & 第2章（情報倫理）、第4章（監査） \\
マネジメント手法 & 第3--4章 & 第1章（経営戦略としてのPDCA） \\
AI/データ活用 & 第7章（AI概論） & 第1章（DX）、第5章（アルゴリズム）
}
\end{hikaku}

\subsection{法律の横断整理}

\begin{hikaku}[法律の出題マップ]
\comptable{法律 & 保護対象/規制内容 & 関連する章}{
個人情報保護法 & 個人情報の適正な取扱い & 第2章・第7章 \\
著作権法 & 創作的な表現 & 第2章 \\
特許法 & 技術的発明 & 第2章 \\
不正アクセス禁止法 & ID/PWの不正使用 & 第2章・第7章 \\
労働者派遣法 & 派遣労働者の保護 & 第2章 \\
特定電子メール法 & 広告メールのオプトイン & 第2章
}
\end{hikaku}

% =============================================================
\section{ひっかけ選択肢の見抜き方}
% =============================================================

\begin{teigi}[ひっかけパターン5選]
\begin{enumerate}
  \item \textbf{似た用語の入れ替え}：SLAとSLM、DNSとDHCP、TCPとUDP——見た目が似た用語を入れ替えて不正解肢に
  \item \textbf{主語のすり替え}：「監査人が改善を実施する」→ 実施は被監査部門。主語に注目
  \item \textbf{「すべて」「必ず」の断定}：「OSSはすべて無料」「暗号化すれば必ず安全」→ 例外がある表現は疑う
  \item \textbf{正しい説明＋別の用語}：「短期間で反復開発→ウォーターフォール」→ 説明は正しいが用語が違う
  \item \textbf{部分的に正しい選択肢}：前半は正しく後半が誤り → 最後まで読むこと
\end{enumerate}
\end{teigi}

\begin{pitfall}[最初の選択肢で安心しない]
\begin{itemize}
  \item 選択肢アが正しそうに見えても、必ずイ〜エまで確認する
  \item 「より適切なもの」を聞かれている場合、アも正しいがイがもっと正確、ということがある
  \item 消去法が有効——明らかに誤りの選択肢を先に消す
\end{itemize}
\end{pitfall}

% =============================================================
\section{本試験の時間配分}
% =============================================================

\begin{teigi}[120分100問の時間配分]
\begin{itemize}
  \item 1問あたり約\textbf{72秒}（120分 ÷ 100問）
  \item \textbf{ストラテジ系}（問1〜35）：約40分
  \item \textbf{マネジメント系}（問36〜55）：約25分
  \item \textbf{テクノロジ系}（問56〜100）：約45分
  \item \textbf{見直し}：10分を確保
\end{itemize}
\end{teigi}

\begin{advice}[時間配分のコツ]
\begin{itemize}
  \item 30秒考えてわからない問題は「仮回答」してマークし、後で戻る
  \item 計算問題は時間がかかるので、知識問題を先に片付ける
  \item 最後の10分で「マーク忘れ」「問題番号のずれ」を確認
\end{itemize}
\end{advice}

\subsection{捨て問の判断基準}

\begin{teigi}[捨て問の判断]
\begin{itemize}
  \item \textbf{捨てるべき問題}：2分以上悩んでも解けない計算問題、見たことのない新用語
  \item \textbf{捨ててはいけない問題}：知識系の問題（消去法で絞れる可能性がある）
  \item 合格基準は600/1000点。100問中60問正解が目安
  \item ただし各分野で300点以上が必要 → 得意分野だけでは合格できない
\end{itemize}
\end{teigi}

% =============================================================
\section{直前チェック10パターン}
% =============================================================

\begin{matome}[試験直前に確認する10パターン]
\begin{enumerate}
  \item \textbf{PLの5段階利益}：売上総→営業→経常→税引前→当期純
  \item \textbf{知的財産権の保護期間}：特許20年、実用新案10年、意匠25年、商標10年（更新可）、著作権死後70年
  \item \textbf{請負/準委任/派遣}：指揮命令権と完成責任の有無
  \item \textbf{V字モデル}：要件定義↔運用テスト、外部設計↔システムテスト
  \item \textbf{MTBF/MTTR/稼働率}：稼働率 $= \text{MTBF} / (\text{MTBF} + \text{MTTR})$
  \item \textbf{TCP vs UDP}：TCP=信頼性、UDP=速度
  \item \textbf{正規化}：第1=繰り返し除去、第2=部分従属除去、第3=推移従属除去
  \item \textbf{公開鍵暗号 vs デジタル署名}：暗号は受信者の公開鍵、署名は送信者の秘密鍵
  \item \textbf{CIA三要素}：機密性・完全性・可用性
  \item \textbf{AI/ML/DL}：AI $\supset$ ML $\supset$ DL の入れ子
\end{enumerate}
\end{matome}

% =============================================================
\section{過去問ドリル}
% =============================================================

\shoumatsuhead{過去問ドリル：横断問題}

\begin{itpmon}[\sourcebadge{R7}\enspace\diffbadge{標}]{%
A社がシステム開発をB社に委託し、B社がC社に再委託している。C社の作業者に対してA社の社員が直接作業指示を行うことは、どのような観点から問題があるか。}
\sentakushi
  {著作権法に違反する可能性がある}
  {偽装請負にあたり、労働者派遣法に違反する可能性がある}
  {個人情報保護法に違反する可能性がある}
  {独占禁止法に違反する可能性がある}
\end{itpmon}
\begin{kaisetsu}{イ}
請負・準委任契約で発注者が直接指揮命令すると\textbf{偽装請負}。法律的には\textbf{労働者派遣法}違反。再委託先への直接指示も同様に問題。第2章の契約形態と第8章の横断視点がつながる。
\end{kaisetsu}

\begin{itpmon}[\sourcebadge{original}\enspace\diffbadge{強}]{%
次の選択肢のうち、記述の内容と用語の組合せが正しいものはどれか。}
\sentakushi
  {データを暗号化して身代金を要求するマルウェア — フィッシング}
  {1〜4週間の短い反復で開発する手法 — ウォーターフォール}
  {サービス品質を数値で取り決めた合意文書 — SLA}
  {故障時に機能を縮退して運転を継続する設計 — フェールセーフ}
\end{itpmon}
\begin{kaisetsu}{ウ}
アはランサムウェアの説明にフィッシングの用語を当てた誤り。イはアジャイルの説明にウォーターフォールを当てた誤り。ウは正しい（SLA=サービスレベル合意）。エはフェールソフトの説明にフェールセーフを当てた誤り。典型的な「説明と用語の入れ替え」ひっかけ。
\end{kaisetsu}

\begin{itpmon}[\sourcebadge{original}\enspace\diffbadge{標}]{%
本試験で残り30分、未回答が20問ある場合の対処として、最も適切なものはどれか。}
\sentakushi
  {すべて空欄のまま見直しに時間を使う}
  {未回答の問題を順に解き、わからない問題も仮回答してマークする}
  {最初の問題から解き直す}
  {難しそうな問題を先に解き、簡単な問題は後回しにする}
\end{itpmon}
\begin{kaisetsu}{イ}
空欄は0点確定。\textbf{仮回答でもマーク}すれば25\%の確率で正解。残り時間が少ないときは「全問マーク」を最優先し、余った時間で見直す。
\end{kaisetsu}

% =============================================================
\section{よくあるミスと到達確認}
% =============================================================

\begin{yokumiss}[この章でよくあるミス]
\begin{enumerate}[leftmargin=1.6em, itemsep=5pt]
  \item \textbf{横断テーマを見落とす}——セキュリティはストラテジ系でも出題される。
  \item \textbf{ひっかけの「すべて」「必ず」に引っかかる}——断定表現は疑う。
  \item \textbf{時間切れで空欄を残す}——空欄は0点。仮回答でもマークすべき。
  \item \textbf{得意分野だけで合格しようとする}——各分野300点以上の足切りがある。
  \item \textbf{計算問題に時間をかけすぎる}——2分以上悩んだら仮回答して先へ進む。
\end{enumerate}
\end{yokumiss}

\begin{tatsaku}
  \item 横断テーマ4つを言え、それぞれの出題例を挙げられる
  \item ひっかけパターン5選をすべて説明できる
  \item 本試験の分野別時間配分を設計できる
  \item 捨て問の判断基準を持っている
  \item 直前チェック10パターンを全部暗唱できる
  \item 各分野の足切りライン（300点）を理解している
  \item 消去法のテクニックを使える
  \item 仮回答→後で見直し の戦略を実行できる
\end{tatsaku}

% =============================================================
% 09_mock.tex — 模擬試験
% =============================================================
\chapter{模擬試験}

\chaptermeta{%
  \item 本試験形式で実力を判定できる
  \item 120分100問のペース配分を体感できる
  \item 弱点分野を特定して復習につなげる
}{180分}{4}{第1--8章すべて}

本章は3セットの模擬試験です。
各セット20問、制限時間30分で解いてください。
\textbf{14問以上正解}を目標にしましょう（本試験の合格ラインに相当）。

\begin{advice}[模擬試験の使い方]
\begin{enumerate}
  \item まず時間を計って通しで解く（答えを見ない）
  \item 採点して正答率を確認する
  \item 間違えた問題は該当する章に戻って復習する
  \item 1週間後にもう一度解き、定着を確認する
\end{enumerate}
\end{advice}

% =============================================================
% セット1：基礎力チェック
% =============================================================

\begin{mogishi}{セット1：基礎力チェック}{30分}{20問中14問}

\begin{itpmon}[\diffbadge{基}]{経営理念の説明として、最も適切なものはどれか。}
\sentakushi
  {経営戦略を現場で実行するための具体的な施策}
  {企業が存在する根本的な目的や価値観を表したもの}
  {市場調査に基づいて策定した販売計画}
  {年度ごとの利益目標を数値化したもの}
\end{itpmon}
\begin{kaisetsu}{イ}
\textbf{経営理念}は企業の存在意義・価値観の表明。アは戦術、ウはマーケティング計画、エは予算。
\end{kaisetsu}

\begin{itpmon}[\diffbadge{基}]{PPMで「市場成長率が高く、市場占有率が低い」事業を何と呼ぶか。}
\sentakushi{花形}{金のなる木}{問題児}{負け犬}
\end{itpmon}
\begin{kaisetsu}{ウ}
成長率高・占有率低 = \textbf{問題児}。投資で花形に育てるか撤退するかを選別する。
\end{kaisetsu}

\begin{itpmon}[\diffbadge{標}]{損益計算書（PL）において「売上総利益 $-$ 販管費」で求められる利益はどれか。}
\sentakushi{売上総利益}{営業利益}{経常利益}{当期純利益}
\end{itpmon}
\begin{kaisetsu}{イ}
売上総利益 $-$ 販管費 $=$ \textbf{営業利益}（本業の儲け）。
\end{kaisetsu}

\begin{itpmon}[\diffbadge{標}]{著作権に関する記述として、誤っているものはどれか。}
\sentakushi
  {著作物を創作した時点で自動的に発生する。}
  {プログラムも著作権で保護される。}
  {アルゴリズムも著作権で保護される。}
  {保護期間は著作者の死後70年である。}
\end{itpmon}
\begin{kaisetsu}{ウ}
\textbf{アルゴリズム}（手順・ロジック）は著作権の対象外。保護されるのは具体的な「表現」。
\end{kaisetsu}

\begin{itpmon}[\diffbadge{基}]{請負契約において、指揮命令権をもつのは誰か。}
\sentakushi
  {発注者}
  {受注者}
  {派遣先企業}
  {発注者と受注者の共同}
\end{itpmon}
\begin{kaisetsu}{イ}
請負契約では\textbf{受注者}が指揮命令権を持つ。発注者が直接指示すると偽装請負。
\end{kaisetsu}

\begin{itpmon}[\diffbadge{標}]{ウォーターフォールモデルの特徴として、最も適切なものはどれか。}
\sentakushi
  {短い反復で段階的に開発する。}
  {前の工程が完了してから次の工程に進み、後戻りしにくい。}
  {試作品を作ってユーザに確認してもらう。}
  {開発チームと運用チームが一体となって作業する。}
\end{itpmon}
\begin{kaisetsu}{イ}
ウォーターフォールは工程を順番に進め、後戻りしにくい。アはアジャイル、ウはプロトタイピング、エはDevOps。
\end{kaisetsu}

\begin{itpmon}[\diffbadge{基}]{V字モデルで、内部設計に対応するテストはどれか。}
\sentakushi{単体テスト}{結合テスト}{システムテスト}{運用テスト}
\end{itpmon}
\begin{kaisetsu}{イ}
内部設計 ↔ \textbf{結合テスト}。プログラミング↔単体、外部設計↔システム、要件定義↔運用。
\end{kaisetsu}

\begin{itpmon}[\diffbadge{標}]{SLAの説明として、最も適切なものはどれか。}
\sentakushi
  {サービスの品質を継続的に改善するプロセス}
  {サービス提供者と利用者の間で合意したサービス品質の取り決め}
  {ITサービスのベストプラクティスをまとめたフレームワーク}
  {利用者からの問い合わせ窓口}
\end{itpmon}
\begin{kaisetsu}{イ}
\textbf{SLA} = Service Level Agreement（サービスレベル合意）。アはSLM、ウはITIL、エはサービスデスク。
\end{kaisetsu}

\begin{itpmon}[\diffbadge{標}]{MTBFが450時間、MTTRが50時間のとき、稼働率はいくらか。}
\sentakushi{0.1}{0.5}{0.9}{0.99}
\end{itpmon}
\begin{kaisetsu}{ウ}
稼働率 $= 450 \div (450 + 50) = 450 \div 500 = 0.9$。
\end{kaisetsu}

\begin{itpmon}[\diffbadge{基}]{インシデント管理の目的として、最も適切なものはどれか。}
\sentakushi
  {インシデントの根本原因を除去する。}
  {サービスの早期復旧を図る。}
  {ITサービスへの変更を管理する。}
  {将来のサービス需要を予測する。}
\end{itpmon}
\begin{kaisetsu}{イ}
\textbf{インシデント管理}は迅速な復旧が目的。アは問題管理、ウは変更管理、エはキャパシティ管理。
\end{kaisetsu}

\begin{itpmon}[\diffbadge{標}]{10進数の7を2進数で表したものはどれか。}
\sentakushi{100}{101}{110}{111}
\end{itpmon}
\begin{kaisetsu}{エ}
$7 = 4 + 2 + 1 = 2^2 + 2^1 + 2^0 = 111_{(2)}$。
\end{kaisetsu}

\begin{itpmon}[\diffbadge{基}]{RAID 1の方式を何というか。}
\sentakushi{ストライピング}{ミラーリング}{パリティ分散}{二重パリティ}
\end{itpmon}
\begin{kaisetsu}{イ}
RAID 1 = \textbf{ミラーリング}。同じデータを2台に書き込む。
\end{kaisetsu}

\begin{itpmon}[\diffbadge{標}]{OSS（オープンソースソフトウェア）の特徴として、最も適切なものはどれか。}
\sentakushi
  {ソースコードが非公開で商用利用が禁止されている。}
  {ソースコードが公開され、利用・改変・再配布が認められている。}
  {メーカーが提供する有償のパッケージソフトウェアである。}
  {特定のOS上でのみ動作するソフトウェアである。}
\end{itpmon}
\begin{kaisetsu}{イ}
\textbf{OSS}はソースコード公開が本質。自由な利用・改変・再配布が認められている。
\end{kaisetsu}

\begin{itpmon}[\diffbadge{標}]{DNSの役割として、最も適切なものはどれか。}
\sentakushi
  {IPアドレスを自動的に割り当てる。}
  {ドメイン名とIPアドレスを対応づける。}
  {電子メールを送信する。}
  {Webページを転送する。}
\end{itpmon}
\begin{kaisetsu}{イ}
\textbf{DNS}はドメイン名 ↔ IPアドレスの変換（名前解決）。アはDHCP、ウはSMTP、エはHTTP。
\end{kaisetsu}

\begin{itpmon}[\diffbadge{標}]{TCPの特徴として、最も適切なものはどれか。}
\sentakushi
  {コネクションレス型で高速な通信を行う。}
  {コネクション型で信頼性の高い通信を行う。}
  {ドメイン名の解決を行う。}
  {IPアドレスを自動割り当てする。}
\end{itpmon}
\begin{kaisetsu}{イ}
\textbf{TCP}はコネクション型で再送制御あり。アはUDP、ウはDNS、エはDHCP。
\end{kaisetsu}

\begin{itpmon}[\diffbadge{基}]{SQLの「SELECT」の役割はどれか。}
\sentakushi
  {テーブルにデータを追加する。}
  {テーブルからデータを取得する。}
  {テーブルのデータを更新する。}
  {テーブルのデータを削除する。}
\end{itpmon}
\begin{kaisetsu}{イ}
\textbf{SELECT}はデータの取得（検索）。アはINSERT、ウはUPDATE、エはDELETE。
\end{kaisetsu}

\begin{itpmon}[\diffbadge{基}]{情報セキュリティの「完全性」が損なわれる事象はどれか。}
\sentakushi
  {許可されていない人がデータにアクセスした。}
  {データが改ざんされた。}
  {システムが停止してサービスが利用できなくなった。}
  {通信内容が盗聴された。}
\end{itpmon}
\begin{kaisetsu}{イ}
\textbf{完全性}（Integrity）はデータの正確性・一貫性。改ざんは完全性の侵害。アとエは機密性、ウは可用性の問題。
\end{kaisetsu}

\begin{itpmon}[\diffbadge{標}]{デジタル署名で確認できることとして、最も適切なものはどれか。}
\sentakushi
  {通信内容の暗号化}
  {送信者の本人確認とデータの改ざん検知}
  {受信者の認証}
  {通信経路の安全性}
\end{itpmon}
\begin{kaisetsu}{イ}
\textbf{デジタル署名}は送信者の\textbf{本人確認}（認証）と\textbf{改ざん検知}（完全性）を実現。暗号化通信は別の仕組み。
\end{kaisetsu}

\begin{itpmon}[\diffbadge{標}]{AI、機械学習、ディープラーニングの関係として正しいものはどれか。}
\sentakushi
  {ディープラーニング $\supset$ 機械学習 $\supset$ AI}
  {機械学習 $\supset$ AI $\supset$ ディープラーニング}
  {AI $\supset$ 機械学習 $\supset$ ディープラーニング}
  {3つは独立した技術分野である}
\end{itpmon}
\begin{kaisetsu}{ウ}
AI $\supset$ ML $\supset$ DL の入れ子構造。AIが最も広い概念。
\end{kaisetsu}

\begin{itpmon}[\diffbadge{基}]{フェールセーフの説明として、最も適切なものはどれか。}
\sentakushi
  {故障しても全体が動き続ける設計}
  {故障時に安全な状態で停止する設計}
  {人間のミスを防ぐ設計}
  {機能を縮退して運転を継続する設計}
\end{itpmon}
\begin{kaisetsu}{イ}
\textbf{フェールセーフ}は故障時に安全側に倒す。アはフォールトトレランス、ウはフールプルーフ、エはフェールソフト。
\end{kaisetsu}

\end{mogishi}

% =============================================================
% セット2：実戦回
% =============================================================

\begin{mogishi}{セット2：実戦回}{30分}{20問中14問}

\begin{itpmon}[\diffbadge{標}]{SWOT分析で扱う4つの要素の組合せとして正しいものはどれか。}
\sentakushi
  {製品・価格・流通・販促}
  {強み・弱み・機会・脅威}
  {顧客・競合・自社・市場}
  {経済価値・希少性・模倣困難性・組織}
\end{itpmon}
\begin{kaisetsu}{イ}
\textbf{SWOT} = Strengths（強み）・Weaknesses（弱み）・Opportunities（機会）・Threats（脅威）。アは4P、エはVRIO。
\end{kaisetsu}

\begin{itpmon}[\diffbadge{標}]{損益分岐点売上高の計算で使う公式はどれか。}
\sentakushi
  {売上高 $\times$ 変動費率}
  {固定費 $\div$ 変動費率}
  {固定費 $\div$ (1 $-$ 変動費率)}
  {(売上高 $-$ 変動費) $\times$ 固定費率}
\end{itpmon}
\begin{kaisetsu}{ウ}
BEP $=$ 固定費 $\div$ ($1 -$ 変動費率) $=$ 固定費 $\div$ 限界利益率。
\end{kaisetsu}

\begin{itpmon}[\diffbadge{標}]{特許権の保護期間として正しいものはどれか。}
\sentakushi
  {登録から10年}
  {出願から10年}
  {出願から20年}
  {著作者の死後70年}
\end{itpmon}
\begin{kaisetsu}{ウ}
\textbf{特許権}は出願から20年。アは商標権（更新可）、イは実用新案権、エは著作権。
\end{kaisetsu}

\begin{itpmon}[\diffbadge{強}]{準委任契約と請負契約の違いとして正しいものはどれか。}
\sentakushi
  {準委任は成果物の完成責任があり、請負にはない。}
  {請負は成果物の完成責任があり、準委任にはない。}
  {準委任では発注者が指揮命令権を持つ。}
  {請負では受注者に善管注意義務がない。}
\end{itpmon}
\begin{kaisetsu}{イ}
\textbf{請負}は成果物の完成責任あり、\textbf{準委任}は善管注意義務のみで完成責任なし。
\end{kaisetsu}

\begin{itpmon}[\diffbadge{標}]{アジャイル開発で使われる「スプリント」の説明として正しいものはどれか。}
\sentakushi
  {プロジェクト全体の工程を一度に計画すること}
  {1〜4週間の短い反復開発の単位}
  {テスト工程だけを繰り返すこと}
  {運用開始後の保守作業の単位}
\end{itpmon}
\begin{kaisetsu}{イ}
\textbf{スプリント}はスクラムにおける1〜4週間のイテレーション。
\end{kaisetsu}

\begin{itpmon}[\diffbadge{標}]{プロジェクトマネジメントのQCDのうち、「D」が表すものはどれか。}
\sentakushi{Design}{Delivery}{Database}{Development}
\end{itpmon}
\begin{kaisetsu}{イ}
QCDのDは\textbf{Delivery}（納期）。Q=Quality（品質）、C=Cost（コスト）。
\end{kaisetsu}

\begin{itpmon}[\diffbadge{標}]{アローダイアグラムのクリティカルパスの説明として正しいものはどれか。}
\sentakushi
  {コストが最も少ない経路}
  {作業数が最も少ない経路}
  {所要時間が最も長い経路}
  {リスクが最も低い経路}
\end{itpmon}
\begin{kaisetsu}{ウ}
\textbf{クリティカルパス}は所要時間が最も長い経路であり、プロジェクト全体の最短所要日数を決定する。
\end{kaisetsu}

\begin{itpmon}[\diffbadge{標}]{問題管理の目的として、最も適切なものはどれか。}
\sentakushi
  {サービスの早期復旧}
  {インシデントの根本原因の除去と再発防止}
  {ITサービスへの変更の管理}
  {SLAの遵守状況の報告}
\end{itpmon}
\begin{kaisetsu}{イ}
\textbf{問題管理}は根本原因の特定と再発防止が目的。アはインシデント管理、ウは変更管理。
\end{kaisetsu}

\begin{itpmon}[\diffbadge{強}]{稼働率0.9のサーバ2台を直列に接続したシステムの稼働率はいくらか。}
\sentakushi{0.81}{0.9}{0.99}{1.8}
\end{itpmon}
\begin{kaisetsu}{ア}
直列の稼働率 $= 0.9 \times 0.9 = \textbf{0.81}$。直列は各稼働率の積。
\end{kaisetsu}

\begin{itpmon}[\diffbadge{標}]{システム監査人の行動として、最も適切なものはどれか。}
\sentakushi
  {被監査部門のシステム改善を直接実施する。}
  {監査結果を被監査部門にのみ報告する。}
  {被監査部門から独立した立場で客観的に評価する。}
  {監査対象のシステム開発にプロジェクトメンバーとして参加する。}
\end{itpmon}
\begin{kaisetsu}{ウ}
監査人は\textbf{独立性}を保ち客観的に評価する。改善の実施は被監査部門の役割。
\end{kaisetsu}

\begin{itpmon}[\diffbadge{基}]{補助記憶装置に該当するものはどれか。}
\sentakushi{レジスタ}{キャッシュメモリ}{メインメモリ}{SSD}
\end{itpmon}
\begin{kaisetsu}{エ}
\textbf{SSD}は補助記憶装置。ア〜ウはいずれも主記憶側（レジスタ→キャッシュ→メインメモリ）。
\end{kaisetsu}

\begin{itpmon}[\diffbadge{標}]{2分探索を行うための前提条件はどれか。}
\sentakushi
  {データが配列に格納されていること}
  {データがソート（整列）済みであること}
  {データの件数が偶数であること}
  {データがリンクリストに格納されていること}
\end{itpmon}
\begin{kaisetsu}{イ}
\textbf{2分探索}の前提は「データがソート済み」であること。未整列では正しく動作しない。
\end{kaisetsu}

\begin{itpmon}[\diffbadge{標}]{DHCPの役割として正しいものはどれか。}
\sentakushi
  {ドメイン名をIPアドレスに変換する。}
  {ネットワーク接続時にIPアドレスを自動的に割り当てる。}
  {電子メールを受信する。}
  {Webページの表示を暗号化する。}
\end{itpmon}
\begin{kaisetsu}{イ}
\textbf{DHCP}はIPアドレスの自動割り当て。アはDNS、ウはPOP3/IMAP、エはHTTPS。
\end{kaisetsu}

\begin{itpmon}[\diffbadge{標}]{RDBにおいて、テーブル間を関連づけるために使われるキーはどれか。}
\sentakushi{主キー}{外部キー}{候補キー}{代替キー}
\end{itpmon}
\begin{kaisetsu}{イ}
\textbf{外部キー}は他テーブルの主キーを参照してテーブル間を関連づける。
\end{kaisetsu}

\begin{itpmon}[\diffbadge{標}]{正規化の第1正規形で排除されるものはどれか。}
\sentakushi
  {繰り返し項目（1セルに複数値）}
  {部分関数従属}
  {推移的関数従属}
  {冗長なインデックス}
\end{itpmon}
\begin{kaisetsu}{ア}
\textbf{第1正規形}は繰り返し項目の排除。イは第2正規形、ウは第3正規形で排除。
\end{kaisetsu}

\begin{itpmon}[\diffbadge{標}]{共通鍵暗号と公開鍵暗号の違いとして正しいものはどれか。}
\sentakushi
  {共通鍵暗号は暗号化と復号で異なる鍵を使い、公開鍵暗号は同じ鍵を使う。}
  {共通鍵暗号は低速で、公開鍵暗号は高速である。}
  {共通鍵暗号は同じ鍵で暗号化と復号を行い、公開鍵暗号はペアの鍵を使う。}
  {共通鍵暗号は鍵の配送が容易で、公開鍵暗号は鍵の配送が難しい。}
\end{itpmon}
\begin{kaisetsu}{ウ}
\textbf{共通鍵}は送受信者が同じ鍵を共有。\textbf{公開鍵}は公開鍵と秘密鍵のペア。アは逆、イも逆、エも逆。
\end{kaisetsu}

\begin{itpmon}[\diffbadge{標}]{標的型攻撃の説明として、最も適切なものはどれか。}
\sentakushi
  {不特定多数に広告メールを大量送信する攻撃}
  {特定の組織や個人を狙い、巧妙な手口で情報を窃取する攻撃}
  {Webサーバに大量のリクエストを送り、サービスを停止させる攻撃}
  {入力フォームに不正なSQL文を挿入する攻撃}
\end{itpmon}
\begin{kaisetsu}{イ}
\textbf{標的型攻撃}は特定の組織や個人を狙った攻撃。アはスパム、ウはDDoS、エはSQLインジェクション。
\end{kaisetsu}

\begin{itpmon}[\diffbadge{標}]{フィッシング攻撃の説明として正しいものはどれか。}
\sentakushi
  {ファイルを暗号化して身代金を要求する。}
  {偽のWebサイトに誘導して個人情報を入力させる。}
  {コンピュータに裏口を仕掛けて遠隔操作する。}
  {キーボード入力を記録して外部に送信する。}
\end{itpmon}
\begin{kaisetsu}{イ}
\textbf{フィッシング}は偽サイトで個人情報を騙し取る。アはランサムウェア、ウはバックドア、エはキーロガー。
\end{kaisetsu}

\begin{itpmon}[\diffbadge{標}]{教師あり学習の例として、最も適切なものはどれか。}
\sentakushi
  {大量のデータから自動的にグループ分けする。}
  {正解ラベル付きのデータでモデルを訓練し、未知データを予測する。}
  {報酬を最大化する行動を試行錯誤で学習する。}
  {人間が作成したルールに基づいて判断する。}
\end{itpmon}
\begin{kaisetsu}{イ}
\textbf{教師あり学習}は正解付きデータで学習。アは教師なし学習（クラスタリング）、ウは強化学習、エはルールベースAI。
\end{kaisetsu}

\begin{itpmon}[\diffbadge{標}]{二要素認証の組合せとして正しいものはどれか。}
\sentakushi
  {パスワードと秘密の質問}
  {パスワードとSMS認証コード}
  {ICカードと指紋認証}
  {イとウの両方が正しい}
\end{itpmon}
\begin{kaisetsu}{エ}
イは知識＋所持、ウは所持＋生体。いずれも異なる2種類の要素を使うため\textbf{二要素認証}。アは知識＋知識で一要素。
\end{kaisetsu}

\end{mogishi}

% =============================================================
% セット3：総仕上げ
% =============================================================

\begin{mogishi}{セット3：総仕上げ}{30分}{20問中14問}

\begin{itpmon}[\diffbadge{標}]{CIOの役職名の正式名称として正しいものはどれか。}
\sentakushi
  {Chief Innovation Officer}
  {Chief Information Officer}
  {Chief Intelligence Officer}
  {Chief Infrastructure Officer}
\end{itpmon}
\begin{kaisetsu}{イ}
\textbf{CIO} = Chief \textbf{Information} Officer（最高情報責任者）。情報システム戦略を統括する。
\end{kaisetsu}

\begin{itpmon}[\diffbadge{強}]{DXの説明として、最も適切なものはどれか。}
\sentakushi
  {MVPで仮説検証を繰り返す開発手法}
  {短期集中でプロトタイプを作り競うイベント}
  {ITを活用してビジネスモデルや企業文化を変革すること}
  {概念の実現可能性を検証する活動}
\end{itpmon}
\begin{kaisetsu}{ウ}
\textbf{DX}（デジタルトランスフォーメーション）はITによるビジネスの変革。アはリーンスタートアップ、イはハッカソン、エはPoC。
\end{kaisetsu}

\begin{itpmon}[\diffbadge{強}]{固定費が800万円、変動費率が0.6のとき、BEPはいくらか。}
\sentakushi{800万円}{1,000万円}{1,333万円}{2,000万円}
\end{itpmon}
\begin{kaisetsu}{エ}
BEP $= 800 \div (1 - 0.6) = 800 \div 0.4 = \textbf{2{,}000}$万円。
\end{kaisetsu}

\begin{itpmon}[\diffbadge{標}]{不正アクセス禁止法で規制される行為はどれか。}
\sentakushi
  {企業の内部告発を行うこと}
  {他人のID/パスワードを使って無断でシステムにログインすること}
  {個人情報を第三者に提供すること}
  {広告メールを無断で送信すること}
\end{itpmon}
\begin{kaisetsu}{イ}
\textbf{不正アクセス禁止法}は他人の認証情報を使った不正ログインを禁止。ウは個人情報保護法、エは特定電子メール法の領域。
\end{kaisetsu}

\begin{itpmon}[\diffbadge{標}]{プロトタイピングモデルの特徴として、最も適切なものはどれか。}
\sentakushi
  {工程を順番に進め、前の工程に戻らない。}
  {短い反復で段階的にシステムを完成させる。}
  {試作品を作ってユーザに確認し、要件を明確にする。}
  {開発と運用を一体化して継続的にサービスを提供する。}
\end{itpmon}
\begin{kaisetsu}{ウ}
\textbf{プロトタイピング}は試作品でユーザの要件を確認する手法。アはウォーターフォール、イはアジャイル、エはDevOps。
\end{kaisetsu}

\begin{itpmon}[\diffbadge{標}]{WBSの目的として、最も適切なものはどれか。}
\sentakushi
  {作業の進捗を横棒グラフで管理する。}
  {プロジェクトの作業を階層的に分解して漏れなく洗い出す。}
  {作業間の依存関係を矢印で表す。}
  {プロジェクトの品質指標を定義する。}
\end{itpmon}
\begin{kaisetsu}{イ}
\textbf{WBS}は作業の階層的分解。アはガントチャート、ウはアローダイアグラム。
\end{kaisetsu}

\begin{itpmon}[\diffbadge{標}]{サービスデスクの別名として使われるものはどれか。}
\sentakushi{SLA}{SPOC}{SLM}{ITIL}
\end{itpmon}
\begin{kaisetsu}{イ}
サービスデスクは\textbf{SPOC}（Single Point of Contact＝単一の問い合わせ窓口）とも呼ばれる。
\end{kaisetsu}

\begin{itpmon}[\diffbadge{強}]{稼働率0.9のサーバ2台を並列に接続したシステムの稼働率はいくらか。}
\sentakushi{0.81}{0.9}{0.99}{1.8}
\end{itpmon}
\begin{kaisetsu}{ウ}
並列の稼働率 $= 1 - (1-0.9)(1-0.9) = 1 - 0.1 \times 0.1 = 1 - 0.01 = \textbf{0.99}$。
\end{kaisetsu}

\begin{itpmon}[\diffbadge{標}]{職務分掌（SoD）の目的として、最も適切なものはどれか。}
\sentakushi
  {業務の効率化を図ること}
  {不正やミスを防ぐために一つの業務を複数の担当者に分離すること}
  {プロジェクトの進捗管理を行うこと}
  {組織の人事評価基準を統一すること}
\end{itpmon}
\begin{kaisetsu}{イ}
\textbf{職務分掌}は業務を分離して一人で完結できない仕組みを作り、不正・ミスを防止する。
\end{kaisetsu}

\begin{itpmon}[\diffbadge{標}]{16進数のAを10進数で表すといくらか。}
\sentakushi{8}{9}{10}{11}
\end{itpmon}
\begin{kaisetsu}{ウ}
16進数のA $=$ 10進数の\textbf{10}。B=11, C=12, D=13, E=14, F=15。
\end{kaisetsu}

\begin{itpmon}[\diffbadge{標}]{XOR（排他的論理和）で「1 XOR 1」の結果はどれか。}
\sentakushi{0}{1}{2}{不定}
\end{itpmon}
\begin{kaisetsu}{ア}
\textbf{XOR}は2つの入力が異なるとき1、同じとき0。$1 \text{ XOR } 1 = 0$。
\end{kaisetsu}

\begin{itpmon}[\diffbadge{標}]{SRAMの主な用途はどれか。}
\sentakushi{メインメモリ}{キャッシュメモリ}{補助記憶装置}{仮想記憶}
\end{itpmon}
\begin{kaisetsu}{イ}
\textbf{SRAM}は高速・高価格でキャッシュメモリに使用。メインメモリにはDRAMを使用。
\end{kaisetsu}

\begin{itpmon}[\diffbadge{標}]{SMTPの役割として正しいものはどれか。}
\sentakushi
  {電子メールの送信}
  {電子メールの受信}
  {Webページの転送}
  {ファイルの転送}
\end{itpmon}
\begin{kaisetsu}{ア}
\textbf{SMTP}はメール送信プロトコル。イはPOP3/IMAP、ウはHTTP、エはFTP。
\end{kaisetsu}

\begin{itpmon}[\diffbadge{強}]{次のSQL文の実行結果として正しいものはどれか。

\texttt{SELECT COUNT(*) FROM 社員表 WHERE 部門 = '営業'}}
\sentakushi
  {営業部門の社員の一覧を表示する。}
  {営業部門の社員の人数を返す。}
  {全部門の社員数を合計して返す。}
  {営業部門の社員の給与合計を返す。}
\end{itpmon}
\begin{kaisetsu}{イ}
\texttt{COUNT(*)}はレコード数を返す集計関数。WHERE句で営業部門に絞っているので、\textbf{営業部門の社員の人数}を返す。
\end{kaisetsu}

\begin{itpmon}[\diffbadge{標}]{サブネットマスク255.255.255.0において、ネットワーク部は何ビットか。}
\sentakushi{8}{16}{24}{32}
\end{itpmon}
\begin{kaisetsu}{ウ}
255.255.255.0 $= 11111111.11111111.11111111.00000000$。上位\textbf{24ビット}がネットワーク部。
\end{kaisetsu}

\begin{itpmon}[\diffbadge{標}]{ソーシャルエンジニアリングの説明として、最も適切なものはどれか。}
\sentakushi
  {ソフトウェアの脆弱性を突いて不正アクセスする手法}
  {人間の心理的な弱みを利用して情報を入手する手法}
  {暗号を解読して通信を傍受する手法}
  {大量の通信でサーバを停止させる手法}
\end{itpmon}
\begin{kaisetsu}{イ}
\textbf{ソーシャルエンジニアリング}は技術的手段ではなく、人間の心理（信頼・恐怖・緊急性）を利用して情報を騙し取る。
\end{kaisetsu}

\begin{itpmon}[\diffbadge{標}]{PKI（公開鍵基盤）において、公開鍵の正当性を証明するのは何か。}
\sentakushi
  {秘密鍵}
  {共通鍵}
  {デジタル証明書（電子証明書）}
  {ハッシュ関数}
\end{itpmon}
\begin{kaisetsu}{ウ}
\textbf{デジタル証明書}は認証局（CA）が発行し、公開鍵の正当性（持ち主の身元）を証明する。
\end{kaisetsu}

\begin{itpmon}[\diffbadge{標}]{過学習（オーバーフィッティング）の説明として正しいものはどれか。}
\sentakushi
  {学習データの偏りによって不公平な判断をすること}
  {訓練データに過剰に適合し、未知データへの予測精度が下がること}
  {AIが人間の知能を超えること}
  {学習に使うデータ量が不足すること}
\end{itpmon}
\begin{kaisetsu}{イ}
\textbf{過学習}は訓練データに適合しすぎて汎化性能が低下する現象。アはAIバイアス、ウはシンギュラリティ。
\end{kaisetsu}

\begin{itpmon}[\diffbadge{強}]{次のうち、横断テーマとして3分野すべてに関連する可能性があるものはどれか。}
\sentakushi
  {損益分岐点の計算}
  {個人情報保護法}
  {RAID}
  {PPM}
\end{itpmon}
\begin{kaisetsu}{イ}
\textbf{個人情報保護法}はストラテジ系（法務）、マネジメント系（監査・内部統制）、テクノロジ系（セキュリティ）の3分野で出題可能性がある横断テーマ。
\end{kaisetsu}

\end{mogishi}

% =============================================================
% 最終チェック
% =============================================================

\begin{tatsaku}[模擬試験の振り返りチェック]
  \item セット1で14問以上正解できた
  \item セット2で14問以上正解できた
  \item セット3で14問以上正解できた
  \item 間違えた問題の該当章を特定できた
  \item 各分野（ストラテジ/マネジメント/テクノロジ）で弱点がないことを確認した
  \item 1問72秒のペースを維持できた
  \item 計算問題でも慌てず公式を適用できた
  \item ひっかけ選択肢を見抜くテクニックを使えた
  \item 消去法で正解を絞り込めた
  \item 本試験に臨む自信がついた
\end{tatsaku}


% =============================================================
% 付録
% =============================================================
\appendix
% =============================================================
% mistakes.tex — 付録A：よくあるミス集
% =============================================================
\chapter*{付録A：よくあるミス集}
\addcontentsline{toc}{chapter}{付録A：よくあるミス集}

全章で学んだ「よくあるミス」を分野横断で整理しました。
試験直前に目を通して、ひっかけパターンを最終確認しましょう。

% =============================================================
% ストラテジ系のミス
% =============================================================
\section*{ストラテジ系のミス}

\begin{yokumiss}[SWOTとVRIOの混同]
SWOT分析は「内部（強み・弱み）×外部（機会・脅威）」の4象限。
VRIO分析は「強みの質」を4つの問い（経済価値・希少性・模倣困難性・組織）で深堀りする。
\textbf{目的が違う}——SWOTは「健康診断」、VRIOは「強みの精密検査」。
\end{yokumiss}

\begin{yokumiss}[PoCとMVPの混同]
\begin{itemize}
  \item \textbf{PoC}（Proof of Concept）= 技術的な実現可能性を確認する小規模な検証
  \item \textbf{MVP}（Minimum Viable Product）= 市場仮説を検証するための最小限の製品
\end{itemize}
PoCは「作れるか？」、MVPは「売れるか？」を確かめる。段階が違う。
\end{yokumiss}

\begin{yokumiss}[PLの5段階利益の順序ミス]
売上総利益→営業利益→経常利益→税引前当期純利益→当期純利益の順。
\begin{itemize}
  \item 営業利益と経常利益を逆にしがち
  \item 「本業の儲け（営業）→ 通常の経営力（経常）」と覚える
  \item 営業外損益（利息、配当など）を含むのが経常利益
\end{itemize}
\end{yokumiss}

\begin{yokumiss}[損益分岐点の公式ミス]
\textbf{正しい公式}：BEP $= \text{固定費} \div (1 - \text{変動費率})$

よくある間違い：
\begin{itemize}
  \item 分母を「変動費率」にしてしまう（逆）
  \item 限界利益率 $= 1 - \text{変動費率}$ を忘れる
  \item 変動費率と固定費率を混同する
\end{itemize}
\end{yokumiss}

\begin{yokumiss}[特許権と著作権の取得方法の混同]
\begin{itemize}
  \item \textbf{特許権}：出願→審査→登録が必要（方式主義）
  \item \textbf{著作権}：創作した時点で自動発生（無方式主義）
\end{itemize}
「著作権は登録が必要」は\textbf{誤り}。特に「プログラムの著作権」と「ソフトウェア特許」は別物。
\end{yokumiss}

\begin{yokumiss}[請負と派遣の指揮命令権]
\begin{itemize}
  \item \textbf{請負}：受注者が指揮命令（発注者は指示しない）
  \item \textbf{準委任}：受注者が指揮命令（請負と同じ）
  \item \textbf{派遣}：派遣先が指揮命令（ここだけ違う）
\end{itemize}
発注者が直接指示 $+$ 請負契約 $=$ \textbf{偽装請負}（違法）。
\end{yokumiss}

% =============================================================
% マネジメント系のミス
% =============================================================
\section*{マネジメント系のミス}

\begin{yokumiss}[V字モデルの対応関係の逆転]
\begin{itemize}
  \item 要件定義 ↔ 運用テスト
  \item 外部設計 ↔ システムテスト
  \item 内部設計 ↔ 結合テスト
  \item プログラミング ↔ 単体テスト
\end{itemize}
「一番上同士、一番下同士」が対応する。外部設計と結合テストを結びつけるのは\textbf{誤り}。
\end{yokumiss}

\begin{yokumiss}[SLAとSLMの逆転]
\begin{itemize}
  \item SLA = Service Level \textbf{Agreement}（合意文書）
  \item SLM = Service Level \textbf{Management}（管理プロセス）
\end{itemize}
「Aは文書、Mは活動」。末尾のアルファベットで区別する。
\end{yokumiss}

\begin{yokumiss}[インシデント管理と問題管理の混同]
\begin{itemize}
  \item インシデント管理 = \textbf{迅速な復旧}（応急処置）
  \item 問題管理 = \textbf{根本原因の除去}（根治治療）
\end{itemize}
「インシデント管理で根本原因を除去する」は\textbf{誤り}。
\end{yokumiss}

\begin{yokumiss}[監査人が「改善を実施する」は誤り]
\begin{itemize}
  \item 監査人の役割：\textbf{助言・勧告}まで
  \item 改善の実施：\textbf{被監査部門}が行う
\end{itemize}
監査人が自ら改善すると\textbf{独立性}が失われる。
\end{yokumiss}

% =============================================================
% テクノロジ系のミス
% =============================================================
\section*{テクノロジ系のミス}

\begin{yokumiss}[2進数変換の余りの読み方]
10進数→2進数の変換で2で割っていく方法。
余りは\textbf{下から上に}読む（最後の余りが最上位ビット）。
上から下に読むと正反対の値になるので注意。
\end{yokumiss}

\begin{yokumiss}[ORとXORの結果の違い]
\begin{itemize}
  \item $1 \text{ OR } 1 = 1$（どちらかが1ならOK）
  \item $1 \text{ XOR } 1 = 0$（同じなら0、違えば1）
\end{itemize}
この違いを問う選択肢が頻出。「排他的」＝「同じならダメ」と覚える。
\end{yokumiss}

\begin{yokumiss}[RAID 0に冗長性があると思い込む]
\begin{itemize}
  \item RAID 0（ストライピング）：\textbf{冗長性なし}。1台故障で全滅
  \item RAID 1（ミラーリング）：2台に同じデータ。1台故障OK
  \item RAID 5（パリティ分散）：3台以上。1台故障まで復旧可能
\end{itemize}
「RAID = 冗長化」ではない。RAID 0はあくまで高速化のため。
\end{yokumiss}

\begin{yokumiss}[公開鍵暗号とデジタル署名の鍵の混同]
\begin{itemize}
  \item \textbf{暗号化通信}：受信者の公開鍵で暗号化 → 受信者の秘密鍵で復号
  \item \textbf{デジタル署名}：送信者の秘密鍵で署名 → 送信者の公開鍵で検証
\end{itemize}
「暗号は受信者の公開鍵、署名は送信者の秘密鍵」——主役が入れ替わる。
\end{yokumiss}

\begin{yokumiss}[フェール○○の混同]
\begin{itemize}
  \item フェール\textbf{セーフ}：故障→\textbf{安全停止}（例：踏切の遮断機）
  \item フェール\textbf{ソフト}：故障→\textbf{機能縮退}で継続（例：エンジン1基停止→残りで飛行）
  \item フォールト\textbf{トレランス}：故障→\textbf{全体は正常動作}（例：RAID）
  \item フール\textbf{プルーフ}：\textbf{人間のミス防止}（例：コンセントの形状）
\end{itemize}
「セーフ＝安全停止」「ソフト＝縮退継続」「トレランス＝耐障害」「プルーフ＝防愚」。
\end{yokumiss}

% =============================================================
% 横断的なミス
% =============================================================
\section*{横断的なミス}

\begin{yokumiss}[「すべて」「必ず」の断定に引っかかる]
選択肢に「すべて」「必ず」「常に」「絶対に」があったら疑う。
\begin{itemize}
  \item 「OSSはすべて無料」→ 誤り（商用サポートは有償可）
  \item 「暗号化すれば必ず安全」→ 誤り（鍵管理が不十分なら意味がない）
  \item 「特許を取れば永久に保護される」→ 誤り（20年の期限あり）
\end{itemize}
\end{yokumiss}

\begin{yokumiss}[似た略語の入れ替え]
試験でよく入れ替えられる組合せ：
\begin{itemize}
  \item SLA ↔ SLM
  \item DNS ↔ DHCP
  \item TCP ↔ UDP
  \item SRAM ↔ DRAM
  \item CIO ↔ CTO
  \item PoC ↔ PoV ↔ MVP
\end{itemize}
それぞれの「一文字の違い」が何を意味するかを押さえておく。
\end{yokumiss}

% =============================================================
% 試験前ミスチェック
% =============================================================

\begin{tatsaku}[試験前ミスチェック]
  \item SWOTとVRIOの目的の違いを言える
  \item PLの5段階利益を順番に言える
  \item BEPの公式を正しく書ける
  \item 特許権（方式主義）と著作権（無方式主義）を区別できる
  \item 請負・準委任・派遣の指揮命令権を即答できる
  \item V字モデルの4組の対応関係を全部言える
  \item SLAとSLMの違いを頭文字で区別できる
  \item インシデント管理（復旧）と問題管理（原因除去）を混同しない
  \item 2進数変換の余りを下から読むことを忘れない
  \item 公開鍵暗号とデジタル署名の鍵の使い分けを即答できる
  \item フェールセーフ/ソフト/トレランス/プルーフを区別できる
  \item 「すべて」「必ず」の断定を疑う習慣がある
\end{tatsaku}

% =============================================================
% reference.tex — 付録B：用語クイックリファレンス
% =============================================================
\chapter*{付録B：用語クイックリファレンス}
\addcontentsline{toc}{chapter}{付録B：用語クイックリファレンス}

試験直前に見返すための用語一覧です。
各分野の頻出用語を1行で整理しています。

% =============================================================
\section*{ストラテジ系の重要用語}
% =============================================================

{
\small
\renewcommand{\arraystretch}{1.3}
\arrayrulecolor{HairLine}
\begin{longtable}{@{}>{\color{Navy}\bfseries\sffamily\footnotesize}p{3.8cm} p{8.2cm}@{}}
\arrayrulecolor{HikakuBar}\toprule
\textbf{用語} & \textbf{一言でいうと} \\
\arrayrulecolor{HairLine}\midrule
\endfirsthead
\arrayrulecolor{HairLine}\midrule
\textbf{用語} & \textbf{一言でいうと} \\
\midrule
\endhead
経営理念 & 企業が存在する根本的な目的・価値観 \\
経営戦略 & 理念を実現するための具体的な道筋 \\
SWOT分析 & 強み・弱み・機会・脅威の4象限で自社を分析 \\
VRIO分析 & 経済価値・希少性・模倣困難性・組織で強みを深堀り \\
PPM & 市場成長率×占有率で事業を4象限に分類 \\
3C分析 & 顧客・競合・自社の3視点で市場を分析 \\
5フォース分析 & 業界の競争圧力を5つの力で分析 \\
4P / 4C & マーケティングミックスの売り手/買い手視点 \\
アンゾフの成長マトリクス & 既存/新規 × 市場/製品で成長方向を4分類 \\
DX & ITでビジネスモデルや企業文化を変革すること \\
PoC & 技術的な実現可能性を確認する小規模な検証 \\
MVP & 市場仮説を検証するための最小限の製品 \\
リーンスタートアップ & MVPで仮説検証を繰り返す手法 \\
損益分岐点（BEP） & 利益がゼロになる売上高 = 固定費 ÷ 限界利益率 \\
売上総利益 & 売上高 $-$ 売上原価（粗利） \\
営業利益 & 売上総利益 $-$ 販管費（本業の儲け） \\
経常利益 & 営業利益 ± 営業外損益 \\
特許権 & 発明を保護（出願から20年） \\
著作権 & 創作的表現を保護（死後70年・無方式主義） \\
偽装請負 & 請負契約で発注者が直接指揮命令する違法状態 \\
\arrayrulecolor{HikakuBar}\bottomrule
\end{longtable}
}

% =============================================================
\section*{マネジメント系の重要用語}
% =============================================================

{
\small
\renewcommand{\arraystretch}{1.3}
\arrayrulecolor{HairLine}
\begin{longtable}{@{}>{\color{Navy}\bfseries\sffamily\footnotesize}p{3.8cm} p{8.2cm}@{}}
\arrayrulecolor{HikakuBar}\toprule
\textbf{用語} & \textbf{一言でいうと} \\
\arrayrulecolor{HairLine}\midrule
\endfirsthead
\arrayrulecolor{HairLine}\midrule
\textbf{用語} & \textbf{一言でいうと} \\
\midrule
\endhead
要件定義 & 「何を作るか」を決める工程（ユーザ主体） \\
外部設計 & 画面・帳票・インターフェースを設計 \\
内部設計 & プログラムの内部構造を設計 \\
ウォーターフォール & 工程を順に進め、後戻りしにくい開発モデル \\
アジャイル開発 & 短い反復で段階的に完成させる開発手法 \\
スクラム & アジャイルの代表的フレームワーク \\
スプリント & スクラムにおける1〜4週間の反復単位 \\
DevOps & 開発と運用の一体化（CI/CD） \\
V字モデル & 設計工程とテスト工程の対応関係 \\
WBS & 作業を階層的に分解した構造 \\
ガントチャート & 横棒で進捗を管理する日程表 \\
クリティカルパス & 最長経路＝プロジェクト全体の最短所要日数 \\
QCD & 品質・コスト・納期の三大制約 \\
SLA & サービス品質の合意文書 \\
SLM & SLAを継続的に管理・改善するプロセス \\
ITIL & ITサービスマネジメントのベストプラクティス集 \\
インシデント管理 & サービスの迅速な復旧（応急処置） \\
問題管理 & インシデントの根本原因を除去し再発防止 \\
MTBF & 平均故障間隔（動いている時間） \\
稼働率 & MTBF ÷ (MTBF + MTTR) \\
システム監査 & 独立した立場でISの信頼性等を評価 \\
職務分掌（SoD） & 業務を分離して不正・ミスを防止 \\
\arrayrulecolor{HikakuBar}\bottomrule
\end{longtable}
}

% =============================================================
\section*{テクノロジ系の重要用語}
% =============================================================

{
\small
\renewcommand{\arraystretch}{1.3}
\arrayrulecolor{HairLine}
\begin{longtable}{@{}>{\color{Navy}\bfseries\sffamily\footnotesize}p{3.8cm} p{8.2cm}@{}}
\arrayrulecolor{HikakuBar}\toprule
\textbf{用語} & \textbf{一言でいうと} \\
\arrayrulecolor{HairLine}\midrule
\endfirsthead
\arrayrulecolor{HairLine}\midrule
\textbf{用語} & \textbf{一言でいうと} \\
\midrule
\endhead
2進数 & 0と1だけで表す数の体系 \\
2の補数 & コンピュータ内部で負の数を表す方式 \\
AND（論理積） & 両方1のときだけ1 \\
OR（論理和） & どちらか一方でも1なら1 \\
XOR（排他的論理和） & 2つが異なるとき1 \\
レジスタ & CPU内部の超高速記憶 \\
キャッシュメモリ & CPUとメインメモリの間の高速バッファ（SRAM） \\
メインメモリ & 実行中プログラムの記憶領域（DRAM） \\
RAID 0 & ストライピング（高速化のみ・冗長性なし） \\
RAID 1 & ミラーリング（同じデータを2台に書込み） \\
RAID 5 & パリティ分散（3台以上・1台故障まで復旧可） \\
OSS & ソースコード公開、自由に利用・改変・再配布可 \\
TCP & コネクション型、信頼性のある通信 \\
UDP & コネクションレス型、高速だが信頼性なし \\
DNS & ドメイン名 ↔ IPアドレスの変換（名前解決） \\
DHCP & IPアドレスの自動割り当て \\
サブネットマスク & ネットワーク部とホスト部の境界を示す \\
主キー & レコードを一意に識別（NULL・重複不可） \\
外部キー & 他テーブルの主キーを参照（テーブル間の関連づけ） \\
SQL（SELECT） & テーブルからデータを取得する命令 \\
正規化 & データの冗長性を排除して更新矛盾を防ぐ \\
CIA三要素 & 機密性・完全性・可用性 \\
公開鍵暗号 & 公開鍵で暗号化、秘密鍵で復号（RSA） \\
デジタル署名 & 送信者の秘密鍵で署名、公開鍵で検証 \\
ランサムウェア & ファイルを暗号化して身代金を要求 \\
フィッシング & 偽サイトで個人情報を騙し取る攻撃 \\
SQLインジェクション & 入力フォームからDB操作する攻撃 \\
フェールセーフ & 故障時に安全な状態で停止する設計 \\
AI / ML / DL & AI $\supset$ ML $\supset$ DL の入れ子構造 \\
AIバイアス & 学習データの偏りによる不公平な判断 \\
\arrayrulecolor{HikakuBar}\bottomrule
\end{longtable}
}

% =============================================================
\section*{直前10パターン}
% =============================================================

\begin{matome}[試験直前に確認する10パターン]
\begin{enumerate}
  \item \textbf{PLの5段階利益}：売上総利益→営業利益→経常利益→税引前当期純利益→当期純利益
  \item \textbf{知的財産権の保護期間}：特許20年、実用新案10年、意匠25年、商標10年（更新可）、著作権 死後70年
  \item \textbf{請負/準委任/派遣}：指揮命令権は請負・準委任＝受注者、派遣＝派遣先
  \item \textbf{V字モデル}：要件定義↔運用テスト、外部設計↔システムテスト、内部設計↔結合テスト、プログラミング↔単体テスト
  \item \textbf{稼働率}：$\text{MTBF} \div (\text{MTBF} + \text{MTTR})$、直列＝積、並列＝$1-(1-A_1)(1-A_2)$
  \item \textbf{TCP vs UDP}：TCP＝信頼性（再送あり）、UDP＝速度（再送なし）
  \item \textbf{正規化}：第1＝繰り返し除去、第2＝部分従属除去、第3＝推移従属除去
  \item \textbf{暗号と署名}：暗号化＝受信者の公開鍵、デジタル署名＝送信者の秘密鍵
  \item \textbf{CIA三要素}：機密性（C）・完全性（I）・可用性（A）
  \item \textbf{AI/ML/DL}：AI $\supset$ ML $\supset$ DL、教師あり/教師なし/強化学習
\end{enumerate}
\end{matome}


\end{document}
